\chapter{Vacuum support and derived actions}
\label{cbc:chap:vacuum-actions}

The scalar vacuum has a nontrivial extension at four positions.
The two components of its support,
$\Delta_{14}$ and $\Delta_{13}\cap\Delta_{24}$,
meet on the total diagonal. A normal coordinate acts
invertibly on the coefficient along the crossed component.
It acts nilpotently on each element of the coefficient at
the intersection. This prevents a splitting of their
extension, even before differential operators are imposed.

Maps of the Heisenberg primitive coefficient must respect
such normal attachments. We first determine its ordinary
horizontal endomorphisms and then its derived maps to the
regular coefficient. The difference between these two
calculations leads to the extension equations for actions.
Those equations retain their prescribed vacuum maps and
the paths comparing compositions.

Throughout, $k$ has characteristic zero. Complexes are
cohomological, with $M[r]^i=M^{i+r}$, and tensor products of
maps use the Koszul rule. All completions will be specified.

\section{The four-position vacuum extension}
\label{cbc:va:four}

Take $X=\mathbb A^1_k$ and use the coordinates
\[
 t=z_4,\qquad a=z_1-z_4,\qquad
 b=z_1-z_3,\qquad c=z_2-z_4.
\]
They are independent. Put $A=k[t,a,b,c]$.
The support of Example~\ref{ex:words-crossed-support} is
\[
 Z_4=V(a)\cup V(b,c),\qquad I(Z_4)=(ab,ac).
\]

For a coordinate $u$, write
$P_u=k[u,u^{-1}]/k[u]$. Its basis consists of $u^{-m}$,
$m\geq1$. Multiplication by $u$ kills $u^{-1}$, and
differentiation sends $u^{-m}$ to $-m u^{-m-1}$.

\begin{proposition}\label{cbc:va:vacuum-extension}
The only nonzero local cohomology modules of $A$ on $Z_4$ are
\begin{align}
 H^1_{Z_4}(A)&=k[t,b,c]\otimes_kP_a,\notag\\
 H^2_{Z_4}(A)&=k[t,a,a^{-1}]\otimes_kP_b\otimes_kP_c.
 \label{cbc:va:four-cohomology}
\end{align}
There is an exact sequence of differential modules
\begin{equation}\label{cbc:va:attachment}
 0\longrightarrow k[t,a]\otimes P_b\otimes P_c
 \longrightarrow H^2_{Z_4}(A)
 \longrightarrow k[t]\otimes P_a\otimes P_b\otimes P_c
 \longrightarrow0.
\end{equation}
It does not split as a sequence of $A$-modules.
\end{proposition}
\begin{proof}
The localization complex for the generators $ab,ac$ is
\[
 A\longrightarrow A[(ab)^{-1}]\oplus A[(ac)^{-1}]
 \longrightarrow A[(abc)^{-1}],
\]
in degrees zero, one and two. Its maps are the diagonal
and the difference of the two inclusions. The first map
is injective. Inside the last localization, the two middle
terms intersect in $A[a^{-1}]$. This gives the first
module of~\eqref{cbc:va:four-cohomology}.

Their sum contains exactly the Laurent monomials with
nonnegative $b$ exponent or nonnegative $c$ exponent.
The remaining monomials have both exponents negative
and arbitrary integral $a$ exponent. They give the
second module. Its nonnegative $a$ powers form the first
term of~\eqref{cbc:va:attachment}. Quotienting leaves
the negative powers of all three normal coordinates.
Multiplication and differentiation preserve this submodule.

Multiplication by $a$ is invertible on the middle term,
whereas each element of the last term is killed by a
power of $a$. Every $A$-linear map from the last term
to the middle term is therefore zero. A section cannot exist.
\end{proof}

The attachment is visible on one class. The indicated
lift of the total-collision class satisfies
\begin{equation}\label{cbc:va:attaching-class}
 a[a^{-1}b^{-1}c^{-1}]=[b^{-1}c^{-1}]\ne0
 \quad\text{in }H^2_{Z_4}(A).
\end{equation}
The right side belongs to the crossed component.
On $a\ne0$, the crossed coefficient remains nonzero
although the outer diagonal is absent.

For the scalar vacuum $J_n=\omega_{X^n}[n]$,
Theorem~\ref{thm:words-support} gives
\[
 (\overline{\Omega_cJ})_4
       \simeq R\Gamma_{Z_4}(\omega_{X^4})[3].
\]
After trivializing the density, its cohomology is
\eqref{cbc:va:four-cohomology} in degrees $-2$ and $-1$.
Thus the local vacuum is already a complex with two
cohomological degrees at four positions. The localizations
and the inverse strong coproducts specify its multiplication.
The two cohomology modules alone do not.

\section{Horizontal symmetries of the primitive coefficient}
\label{cbc:va:rigidity}

Return to the left differential modules of
Chapter~\ref{cbc:chap:heisenberg}. Retain
$R_n=k[z_1,\ldots,z_n]$, its Weyl algebra $D_n$, the
primitive module $H_n$, and its generators $g_S$.
The coefficient of $w^{-1}$ in the expansion at infinity
of a proper rational representative defines total finite residue:
\[
 \operatorname{res}:H_n\longrightarrow R_n,
 \qquad \operatorname{res}(g_i)=1,\qquad
 \operatorname{res}(g_S)=0\quad(|S|>1).
\]
This coefficient lies in $R_n$, and a $w$-derivative
has zero residue. Hence the map is well defined and horizontal.

For $v=(v_1,\ldots,v_n)\in k^n$, put
\[
 s_v:R_n\longrightarrow H_n,\qquad
 s_v(1)=\sum_i v_i g_i.
\]
Each $g_i$ is horizontal, so $s_v$ is a map of $D_n$-modules.

\begin{theorem}\label{cbc:va:primitive-endomorphisms}
For $n\geq2$, every $D_n$-linear endomorphism of $H_n$
has a unique expression
\begin{equation}\label{cbc:va:endomorphism-form}
 T=\lambda\,\mathrm{id}+s_v\operatorname{res},
 \qquad \lambda\in k,\quad v\in k^n.
\end{equation}
In residue coordinates its matrix is
$\lambda I+v\mathbf1^{\mathsf t}$. Composition is
\[
 (\lambda,v)(\mu,w)
  =(\lambda\mu,\lambda w+\mu v+
                       (\mathbf1^{\mathsf t}w)v).
\]
The scalar target and the residue kernel satisfy
\[
 \operatorname{Hom}_{D_n}(H_n,R_n)=k\operatorname{res},
 \qquad \operatorname{End}_{D_n}(\ker\operatorname{res})=k.
\]
In particular, the residue kernel is indecomposable.
\end{theorem}
\begin{proof}
Flatness embeds $H_n$ in its generic residue space
$k(z_1,\ldots,z_n)^n$. The connection there is ordinary
position differentiation. A horizontal matrix therefore
has entries in $k$.

Near the generic point of $z_i=z_j$, all other differences
are invertible. The binary calculation
\eqref{cbc:hc:binary-transition} gives one localized
direction $e_i-e_j$ and regular complementary directions.
An element in the image of a linear map from the localized
direction to a regular direction would be divisible by
every power of $z_i-z_j$. The regular module has no such
nonzero element. The matrix must therefore preserve
$k(e_i-e_j)$.

For $n\geq3$, the relation
$(e_i-e_j)+(e_j-e_k)=e_i-e_k$ forces the eigenvalues on
these lines to agree. For $n=2$ there is one line.
Thus, in either case, $T-\lambda I$ kills the sum-zero
hyperplane and has the form $v\mathbf1^{\mathsf t}$.
The maps in~\eqref{cbc:va:endomorphism-form} realize all
these matrices on $H_n$, proving the assertion and its
composition law.

A horizontal map $H_n\to R_n$ is generically a constant
row vector. Regularity at every binary diagonal forces
it to kill every $e_i-e_j$. It is a multiple of residue.

The map $s_{e_1}$ splits residue as a differential-module
map, giving $H_n=R_n\oplus\ker\operatorname{res}$.
Extend an endomorphism of the kernel by zero on this
regular summand. Formula~\eqref{cbc:va:endomorphism-form}
restricts to $\lambda I$ on the kernel. Its endomorphism
ring is $k$, which has no nontrivial idempotents.
\end{proof}

The splitting by $s_{e_1}$ singles out a position.
Compatibility with the whole family of coefficients
removes this choice.

\begin{proposition}\label{cbc:va:natural-scalars}
A family of horizontal endomorphisms $T_n$ of $H_n$
is multiplication by one scalar independent of $n$
under either of the following hypotheses:
\begin{enumerate}
 \item it commutes with every permutation and every
       coefficient collision map of
       Proposition~\ref{cbc:hc:collisions}.
 \item it commutes with the block decompositions of
       Proposition~\ref{cbc:hc:cuts}. One nontrivial
       consecutive cut in each arity $n\geq2$ suffices.
\end{enumerate}
\end{proposition}
\begin{proof}
Under the first hypothesis, permutation compatibility
gives $T_n=\lambda_nI+c_n\mathbf1\mathbf1^{\mathsf t}$
for $n\geq2$. On the simple-pole classes, a collision
$\alpha$ sends $g_i$ to $g_{\alpha(i)}$.
A three-to-two collision of fiber sizes $(2,1)$ therefore
gives, for both possible output labels,
\[
 \lambda_3e_j+c_3(2e_1+e_2)
       =\lambda_2e_j+c_2(e_1+e_2).
\]
Subtract the equations for $j=1,2$ and then compare
their coordinates. One obtains
$\lambda_3=\lambda_2$ and $2c_3=c_2=c_3$, hence
$c_2=c_3=0$.

For $n>3$, a collision with fiber sizes
$(n-1,1)$ similarly gives $c_n=0$ and
$\lambda_n=\lambda_2$. Collision to a singleton fixes
the last scalar.

For the second hypothesis, write
$T_n=\lambda_nI+v_n\mathbf1^{\mathsf t}$.
A cut into nonempty blocks of sizes $p,q$ forces its
off-diagonal matrix blocks to vanish. They are
\[
 v_{n,[1,p]}\mathbf1_q^{\mathsf t},
 \qquad v_{n,[p+1,n]}\mathbf1_p^{\mathsf t}.
\]
Thus $v_n=0$, and the two diagonal blocks give
$\lambda_n=\lambda_p=\lambda_q$. Induction from
$\operatorname{End}_{D_1}(H_1)=k$ proves the result.
\end{proof}

On the diagram with monotone collisions alone, the maps
$g_1\operatorname{res}$ and $g_n\operatorname{res}$ are
natural and distinct for $n\geq2$. A monotone surjection
preserves the first and last labels, and total residue
commutes with substitution. They fail the cut condition:
on the cut $[1,2]\mid[3,4]$, the first map sends the
residue vector $e_3$ to $e_1$. Thus the cut maps supply
an essential part of the rigidity assertion.

\section{Derived scalar maps and the residue simplex}
\label{cbc:va:derived-residue}

Ordinary scalarity does not determine derived endomorphisms.
The simplest comparison uses the regular module as target.
We compute it from the normal directions of the diagonal
constituents, retaining their nonsplit attachments.

For an affine linear subspace $L\subseteq\mathbb A^n$,
let $\delta_L$ denote the trivial left differential
module on $L$, pushed forward with all normal derivatives.

\begin{lemma}\label{cbc:va:normal-hom}
Let affine linear subspaces $L,M\subseteq\mathbb A^n$
have nonempty intersection. If their dimensions are
$l,m$, and $d=\dim(L\cap M)$, then choices of linear
normal coordinates give
\[
 R\operatorname{Hom}_{D_n}(\delta_L,\delta_M)
              \simeq k[-(l+m-2d)].
\]
\end{lemma}
\begin{proof}
Translate a common point to zero and split the ambient
space into common tangents, tangents of $L$ alone,
tangents of $M$ alone, and common normals.
In one variable put $O=k[u]$ and retain $P_u$ above.
Their left Weyl-module resolutions use right
multiplication by $\partial_u$ and by $u$, respectively.
Applying Hom gives the two-term complexes
\[
\begin{array}{c|c|c}
 (\text{source},\text{target})&d&\text{cohomology representative}\\\hline
 (O,O)&\partial_u:O\to O&1\text{ in degree }0\\
 (P_u,P_u)&u:P_u\to P_u&u^{-1}\text{ in degree }0\\
 (O,P_u)&\partial_u:P_u\to P_u&u^{-1}\text{ in degree }1\\
 (P_u,O)&u:O\to O&1\text{ in degree }1.
\end{array}
\]

Both terms have degrees zero and one. Right inverses
on the boundaries are, in the displayed order,
\[
 u^j\mapsto\frac{u^{j+1}}{j+1},\quad
 u^{-j}\mapsto u^{-j-1},\quad
 u^{-j}\mapsto-\frac{u^{-j+1}}{j-1}\ (j\geq2),\quad
 u^j\mapsto u^{j-1}\ (j\geq1).
\]
Set these homotopies to zero on the degree-one
cohomology representatives. They contract the complementary
monomials. Tensor their contractions with the Koszul signs.
Exactly the $l+m-2d$ directions tangent to only one
subspace contribute degree one.
\end{proof}

\begin{theorem}\label{cbc:va:scalar-simplex}
The derived scalar mapping complex has a model
\begin{equation}\label{cbc:va:simplex}
 R\operatorname{Hom}_{D_n}(H_n,R_n)
 \simeq
 \left[k^{\binom n1}\longrightarrow k^{\binom n2}
       \longrightarrow\cdots\longrightarrow k^{\binom nn}\right]
 \simeq k.
\end{equation}
Its terms have degrees zero through $n-1$. With subsets
in increasing order, the differential is
\[
 (dc)_{i_0\cdots i_p}
   =\sum_{j=0}^p(-1)^j
             c_{i_0\cdots\widehat{i_j}\cdots i_p}.
\]
The surviving map is total finite residue.
\end{theorem}
\begin{proof}
Apply derived Hom into $R_n$ to the finite filtration
of Theorem~\ref{cbc:hc:filtration}. A constituent
indexed by a subset of size $s$ has mapping cohomology
$k$ in degree $s-1$, by Lemma~\ref{cbc:va:normal-hom}.
Choose the normal contractions in its proof and transfer
along this finite filtration. A transferred differential
has degree one, so it joins only adjacent subset sizes.

Let $T=S\setminus\{i\}$, and choose $r\in T$.
The attaching relation
\[
 (z_i-z_r)g_S=g_{S\setminus\{r\}}-g_T
\]
has coefficient $-1$ on $g_T$. In that component its
new normal direction is the localization extension
$0\to O\to O[u^{-1}]\to P_u\to0$.
The resulting incidence entry is therefore a unit.

For $|S|\geq3$, a subset of size $|S|-1$ not contained
in $S$ has no degree-one extension with its diagonal,
by Lemma~\ref{cbc:va:normal-hom}.
At size two all singleton constituents have the same
support. The actual relation
$(z_i-z_j)g_{ij}=g_i-g_j$ fixes the two opposite
endpoint entries and makes every spectator entry zero.

Order the normal coordinates to fix their orientations.
The edge entries are the displayed differences.
Square-zero on each two-step face fixes the higher
incidence signs after one choice of orientation per subset.
This gives the complex~\eqref{cbc:va:simplex}.

Choose vertex $1$. In degree $p\geq1$ define
\[
 (hc)_S=
 \begin{cases}
 c_{\{1\}\cup S},&1\notin S,\\
 0,&1\in S,
 \end{cases}
\]
with $1$ first in the orientation. In degree zero put $h=0$.

Let $\pi(c)=c_{\{1\}}$ in degree zero and let $i$ include
constant vertex cochains. Cancellation of incidences gives
$dh+hd=1-i\pi$. Thus only the common value on all $g_i$
survives, which is total residue. The finite transfer
lifts this contraction to the derived mapping complex.
\end{proof}

For $n=3$ the two differentials are
\[
 (c_1,c_2,c_3)\longmapsto(c_2-c_1,c_3-c_1,c_3-c_2),
 \qquad
 (x_{12},x_{13},x_{23})\longmapsto x_{23}-x_{13}+x_{12}.
\]
The choice of vertex specifies the contraction, not
the residue map, which is independent of that choice.

The endomorphism target gives a different answer already
at two positions. Suppress common translation, put
$R=k[u]$, $L=R[u^{-1}]$, and use
$H_2\simeq R\oplus L$ from~\eqref{cbc:hc:binary-transition}.
Proposition~\ref{cbc:hc:reverse-hom} gives
$R\operatorname{Hom}_D(L,R)=0$ and
$\operatorname{Ext}^0_D(L,L)=
\operatorname{Ext}^1_D(L,L)=k$.

The other two corners are $[R\xrightarrow{\partial_u}R]$
and $[L\xrightarrow{\partial_u}L]$.
The first has only its constant class. The second has
a constant class in degree zero and the class of $u^{-1}$
in degree one.

Consequently
\[
 \dim_k\operatorname{End}_D(H_2)=3,\qquad
 \operatorname{Ext}^1_D(H_2,H_2)=k^2,\qquad
 \operatorname{Ext}^{q}_D(H_2,H_2)=0\quad(q>1).
\]
Tensoring back the common tangential regular module
does not change these groups. To compute derived
natural endomorphisms of the whole coefficient system,
one must impose the collision and cut comparisons with
their homotopies. Neither the single-arity endomorphism
calculation nor~\eqref{cbc:va:simplex} performs that step.

\section{Pointed actions and their comparison paths}
\label{cbc:va:pointed}

A vacuum-preserving map of coefficients is only the
linear part of an action. To compare actions, one must
also retain the path that identifies the restricted
action with the given one.

The action calculations use a symmetric convolution
on the full ordered collision atlas. We specify this
carrier before forming its mapping spaces.
On the affine line, its underlying scheme is the
ravioli configuration space of Alfonsi--Kim--Young
\cite{AKY2024}*{\S4.1}.

For a nonempty finite set $I$, take one copy $X^I_\sigma$
of $X^I$ for each total order $\sigma$ of $I$.
For two orders $\sigma,\rho$, remove the diagonal
$z_i=z_j$ whenever they disagree on the order of $i,j$,
and identify the resulting opens by their position
coordinates. Denote the glued scheme by $\mathscr R_I$
and its position map by $q_I:\mathscr R_I\to X^I$.

These are valid gluing data. If two orders disagree
on a pair, at least one disagrees with any third
order on that pair. Thus
$V_{\sigma\rho}\cap V_{\sigma\eta}
=V_{\sigma\rho}\cap V_{\rho\eta}$ for the prescribed
overlaps. The finite smooth chart cover makes
$\mathscr R_I$ smooth, quasi-compact and quasi-separated,
and $q_I$ is locally an isomorphism.
At a position tuple with equality blocks $B$, its
fiber has $\prod_B|B|!$ points, one for each choice
of order inside every equality block.
In particular the total triple collision has six
branches.

A surjection $\alpha:I\twoheadrightarrow H$ with
ordered fibers expands an order on $H$ by concatenating
those fibers. The diagonal substitution then defines
$\delta_\alpha:\mathscr R_H\to\mathscr R_I$.
Indeed the preimage of an overlap between expanded
orders is exactly the overlap of the original orders.
Concatenation of fiber orders proves
\mbox{$\delta_\alpha\delta_\beta=\delta_{\beta\alpha}$}.
For example, a collision identifying $1,3$ on the
chart $123$ reaches that chart only where $z_1=z_3\ne z_2$.
Its selected total-collision branches occur on charts
where $1,3$ are consecutive.

Let $\mathscr D$ be the enhanced category of right
differential-module families on these schemes with
coherent equivalences
$\delta_\alpha^!M_I\simeq M_H$, together with the
separate empty coefficient in $D(k)$.
Differential modules on $\mathscr R_I$ are defined by
descent on its finite smooth chart cover.
Use the derived six-functor operations with the
cochain and density conventions of
Chapter~\ref{chap:collision-words}.

For $I=A\amalg B$, let $U_{A,B}$ separate all cross
pairs, and let $j_{A,B}$ be its open inclusion.
Restriction of orders identifies this open with
$(\mathscr R_A\times\mathscr R_B)_{\mathrm{disj}}$:
all shuffles of the two block orders agree there.
Define
\begin{equation}\label{cbc:va:symmetric-convolution}
 (M\star N)_I=
 \bigoplus_{I=A\amalg B}
 j_{A,B,*}(M_A\boxtimes N_B)|_{U_{A,B}}.
\end{equation}
The pairs of subsets are ordered. Either may be empty.

As in Proposition~\ref{prop:words-cut-monoidal},
all iterated products have the common coefficient
on the open separating their blocks. This proves
associativity and its higher coherences. Permuting
blocks gives the symmetry.

Extraordinary pullback
kills a summand whose two groups meet inside one
collision fiber. The other summands are exactly
the decompositions of the quotient. Base change
therefore gives the required collision maps.
The unit is $K$, supported at the empty configuration.
Every product of more than $|I|$ positive factors
vanishes on $\mathscr R_I$.

Restricting to the increasing chart and projecting
onto consecutive decompositions gives the word-cut
comparison. No equivalence of the two diagram
categories is asserted by this projection.
In particular their derived mapping spaces must
be distinguished.

Now let $\mathscr C$ be the full subcategory of
coaugmented counital coassociative coalgebras for
$\star$ whose coproducts are equivalences on every
block-disjoint open. It carries the diagonal tensor
$\otimes^!$, with unit the local vacuum $J$.
Its coproduct uses the pairs of equal decompositions
in the two finite sums~\eqref{cbc:va:symmetric-convolution}.
For those summands the two localizations have the
same open and combine by base change.
Iterated equal decompositions verify coassociativity
and the tensor compatibilities.

For a coalgebra $F$ with a pointing $u:J\to F$, and
$C\in\mathscr C$, define its space of pointed actions by
\begin{equation}\label{cbc:va:pointed-action}
 \Phi_C(F,u)=
 \operatorname{fib}_{\mathrm{id}_C}
 \left(
 \operatorname{Map}_{\mathscr C}(F\otimes^!C,C)
 \longrightarrow\operatorname{Map}_{\mathscr C}(C,C)
 \right).
\end{equation}
The arrow restricts along $u\otimes^!\mathrm{id}_C$.
An element therefore includes both a map $\alpha$ and
a specified path from its restriction to $\mathrm{id}_C$.
All mapping spaces are derived.

A morphism from $(F,u,\alpha)$ to $(G,v,\beta)$ consists
of a pointed map $f:F\to G$ and a path
$\alpha\simeq\beta(f\otimes^!\mathrm{id}_C)$ compatible
with the two pointings. Its higher simplices retain
the corresponding higher paths. Denote this category
by $\mathscr A_C$.
A representing pointed coalgebra $Z$ would give
natural equivalences
\[
 \operatorname{Map}_{J/}(F,Z)\simeq\Phi_C(F,u).
\]
The equivalences must hold for every permitted actor $F$.

There is a second question when $F$ is a transverse
monoid, meaning an associative algebra for $\otimes^!$.
An action of this monoid includes its multiplication
and unit coherences. Let $\operatorname{Act}_F(C)$
be the space of such coherent actions.

A \emph{local center} is a final object among transverse
monoids acting on $C$, with $C$ fixed. For a candidate
$(Z,\beta)$ and a test action $(F,\alpha)$, its mapping
space is
\begin{equation}\label{cbc:va:action-fibre}
 \operatorname{fib}_{\alpha}
 \left(
 \operatorname{Map}_{\operatorname{Alg}(\mathscr C)}(F,Z)
 \longrightarrow\operatorname{Act}_F(C)
 \right).
\end{equation}
The map composes with $\beta$. Thus a loop of test
actions outside the image of the source fundamental
group makes this fiber disconnected. It disproves
finality of that candidate.

\subsection{Free fields and the diagonal tensor}

We will use the field coalgebras of
Chapter~\ref{cbc:chap:transfer} on the affine line.
A free-field vertex algebra here has finitely many
homogeneous generators, scalar polynomial
$\lambda$-brackets, and a basis of ordered monomials
in their translates. Its differential is linear in
those translates. The examples below are tensor
products of Heisenberg fields, regular commutative
fields and Weyl--Clifford fields.

Let $L_B$ be the positive diagonal family of the
state differential module of $B$.
For $|I|\geq2$ its coefficient is the sum of the
full normal direct images along the small diagonal
in every order chart. These diagonals have disjoint
images: two different order charts overlap only
away from at least one of their pair diagonals.
The singleton coefficient is the state differential
module itself. The empty coefficient is zero.
The operator-product residues define a Lie bracket
$L_B\star L_B\to L_B$.

Its Chevalley coalgebra is
\[
 \operatorname{CE}_\star(L_B)
   =K\oplus\bigoplus_{r\geq1}
                      \operatorname{Sym}^r_\star(L_B[1]).
\]
It has the unshuffle coproduct, and its
coderivation is determined by
\[
 Q_1(sx)=-s(dx),\qquad
 Q_2(sx\,sy)=(-1)^{|x|}s[x,y].
\]
The square-zero identities are respectively the
internal chain equation, the derivation rule and
Jacobi. Positivity makes the displayed sum finite
at each configuration.

For these algebras $B$, denote this full residue
coalgebra by $C_B$. In its partition description,
each summand is the normal direct image from the
diagonals of its blocks, localized away from
collisions between blocks. $Q_2$ is their residue
boundary. Thus it is also the partition Cousin
model of the factorization lattice on the order
charts. Its coproduct keeps all block decompositions.

The consecutive components give the cut structure
used earlier. The vacuum map $k\to B$ gives
$J=C_k\to C_B$.

\begin{lemma}\label{cbc:va:chevalley-bar}
For a positive dg Lie object $L$, the twisting
cochain $sx\mapsto x$ induces an equivalence
\[
 \operatorname{CE}_\star(L)\xrightarrow{\sim}B_\star U_\star L.
\]
Its length-$r$ component is the unnormalized
suspended symmetrization
\[
 sx_1\cdots sx_r\longmapsto
 \sum_{\sigma\in S_r}\epsilon(\sigma;sx_1,\ldots,sx_r)
       [x_{\sigma(1)}|\cdots|x_{\sigma(r)}],
\]
where $\epsilon$ is the Koszul permutation sign.
\end{lemma}
\begin{proof}
Deconcatenating the permutation sum chooses a
subset and independently orders it and its
complement. This is its unshuffle coproduct.
The binary chain equation is the envelope
relation
$xy-(-1)^{|x||y|}yx=[x,y]$. Grouping adjacent
pairs proves it in every length.

Filter the envelope by the number of original
$L$ inputs. The multilinear PBW identity gives
$\operatorname{gr}U_\star L=\operatorname{Sym}_\star L$:
ordering a triple has the same result precisely
by Jacobi, and averaging is exact in characteristic
zero. Apply this identity to a free simplicial
Lie resolution for derived inputs.
The associated graded comparison is the abelian
Koszul comparison.

Its free resolution is
$S\star\operatorname{Sym}_\star(L[1])$, where
$S=\operatorname{Sym}_\star L$ and
$d_K(sx)=-x$. The derivation $h(x)=-sx$,
$h(sx)=0$ satisfies $d_Kh+hd_K=N$, with
$N$ the total polynomial weight.
It anticommutes with the internal differential.
Thus $h/N$ contracts each positive weight.
The permutation sum compares this resolution
with the free bar resolution and induces the
identity on their common augmentation.
Tensoring with the augmentation module proves
the associated graded equivalence.

Both filtrations stop at $|I|$ on a configuration
$I$, so filtered comparison proves the lemma.
The coalgebra map has no factor $1/r!$.
That factor belongs to PBW averaging.
\end{proof}

\begin{lemma}\label{cbc:va:free-field-tensor}
For free fields $A,B$, the residue construction has
natural equivalences
\[
 C_A\otimes^!C_B\simeq C_{A\otimes B}.
\]
They preserve units, collisions and iterated
factorization. The construction preserves
quasi-isomorphisms of dg vertex algebras in this
class and polynomial homotopies between their maps.
\end{lemma}
\begin{proof}
The free-field factorization lattice on a position
chart is the vacuum module for polar sections at the
marked points
\cite{BD04}*{\S\S3.7.7, 3.7.20}.
Filter this module by the number of generating fields.
Its associated graded is the graded symmetric algebra
of their principal parts. Those principal parts are
free over the position ring by~\eqref{cbc:hc:free-basis}.

At a fixed polynomial degree the complex is bounded
and flat, because there are finitely many generating
cohomological degrees. The differential preserves the
filtration. Successive extensions and filtered unions
therefore give a $K$-flat lattice. Here $K$-flat means
that tensoring with it preserves acyclic complexes.

The fields of $A$ have zero bracket with those of $B$.
Their Lie algebras of polar sections consequently
have commuting noncentral factors. Their joint vacuum
representation imposes each scalar central relation at
its stated value and is the tensor product of their
vacuum modules. This identifies the lattice of
$A\otimes B$ with the tensor product of the lattices,
compatibly with collisions and separation of groups.

The finite Cousin resolution of a lattice is its
partition residue complex: the normal boundary maps
are the oriented residues of
Section~\ref{cbc:ct:residues}. Tensoring $K$-flat lattices
computes the derived tensor. Passing to right
differential modules contracts one copy of the density
unit $J$, giving the asserted $\otimes^!$ comparison.
The maps come from polar sections, so their
associativity, symmetry, unit and collision diagrams
commute before passing to the derived category.

A quasi-isomorphism gives a quasi-isomorphism on each
derived partition coefficient. The finite partition
filtration then proves the assertion with its residue
differential. For a polynomial homotopy use the fixed
target tensored with polynomial forms. Evaluation and
constant inclusion are inverse quasi-isomorphisms of
the parameter complex. Their maps of residue
coalgebras retain all faces.
\end{proof}

The same argument for a free regular commutative field
uses its multijet lattice
\cite{BD04}*{\S\S3.4.20--3.4.22}.
If $P_I(w)=\prod_{i\in I}(w-z_i)$, each jet stage
$R_I[w]/P_I^m$ is finite free. Chinese remainders
give separation of groups, and cofinal powers give
collision. The positive cohomological degree of a
generating field changes its graded symmetric powers
but not this construction.

\subsection{A cochain action}
\label{cbc:va:cochain-action}

Fix $\kappa\in k^\times$ and let $V_\kappa$ be the
Heisenberg vertex algebra, with even current $J$ and
$[J_\lambda J]=\kappa\lambda$.
Adjoin a regular odd field $c$ of degree one and put
\[
 \mathcal Z_\kappa=V_\kappa\otimes\Lambda[c,Tc,T^2c,\ldots],
 \qquad dJ=\kappa Tc,\quad dc=0.
\]
The differential is the zero mode $(:cJ:)_{(0)}$.
It is a derivation of every vertex product and
squares to zero.

Introduce even fields $t,p$ of degree zero and odd
fields $\eta,\iota$ of degrees $1,-1$, with brackets
\[
 [p_\lambda t]=1,\qquad [\iota_\lambda\eta]=1.
\]
All other generating brackets vanish, apart from
graded skewsymmetry. Their quartet algebra $E$ has
differential
\[
 Qt=\eta,\qquad Q\iota=p,\qquad Q\eta=Qp=0.
\]

On $W=V_\kappa\otimes E$, let $Q$ act trivially on
the first factor. The assignments
\begin{equation}\label{cbc:va:quartet-maps}
 i(J)=J+p,\qquad
 j(J)=J-\kappa Tt,\qquad j(c)=-\eta
\end{equation}
define dg vertex maps $i:V_\kappa\to W$ and
$j:\mathcal Z_\kappa\to W$.
Their images commute: for example,
$[J+p{}_\lambda J-\kappa Tt]=\kappa\lambda-\kappa\lambda=0$.
Thus
\[
 \rho:\mathcal Z_\kappa\otimes V_\kappa\longrightarrow W,
 \qquad \rho(a\otimes v)=j(a)_{(-1)}i(v)
\]
is a dg vertex map and restricts to $i$ on the vacuum axis.

The map $i$ is a quasi-isomorphism. In the ordered
creation-mode basis use $J+p$ in place of $J$.
The remaining variables form the contractible pairs
$(T^nt,T^n\eta)$ and $(T^n\iota,T^np)$.
Their Euler contraction sends $\eta$ to $t$ and $p$
to $\iota$, with the corresponding formulas for
translates, and divides by positive total quartet
degree. Each state is a finite polynomial, so the
contraction is finite. It commutes with translation.

Give $J,p,\iota$ conformal weight one and $t,\eta$
weight zero, and let $T$ raise weight by one.
Both the differential and this contraction preserve
conformal weight. In particular, they preserve the
subspace of positive conformal weight.

Lemma~\ref{cbc:va:free-field-tensor} now defines a pointed
action $\alpha_{\mathcal Z}$ by the roof
\begin{equation}\label{cbc:va:cochain-roof}
 C_{\mathcal Z_\kappa}\otimes^!C_{V_\kappa}
 \simeq C_{\mathcal Z_\kappa\otimes V_\kappa}
 \xrightarrow{C_\rho}C_W
 \xleftarrow[\sim]{C_i}C_{V_\kappa}.
\end{equation}
The inverse of $C_i$ is taken in the derived category.
The vacuum axis supplies its specified pointing.
The obstruction below depends on this action. It
applies to every coherent transverse monoid structure
on $C_{\mathcal Z_\kappa}$ with this underlying action.

\section{A vacuum probe against finality}
\label{cbc:va:nonfinality}

Let $C=C_{V_\kappa}$ and let
\[
 N=k[t_0,t_1,\ldots],\qquad |t_n|=2,\qquad
 Tt_n=t_{n+1},\qquad dN=0.
\]
This is the free regular commutative vertex algebra
on $t_0$. Write $\epsilon:N\to k$ for its augmentation
and $F=C_N$. Multiplication in $N$ makes $F$ a
commutative transverse monoid. Its augmentation
defines a coherent trivial action $\alpha_0$ on $C$.
The composite $g_0:F\to J\to C_{\mathcal Z_\kappa}$
is a pointed map inducing that action.

We use total polynomial degree in the $t_n$, denoted
by $G$, to separate the deformation equations.
This is independent of their cohomological degree $2G$.
Put $\mathfrak g=L_{V_\kappa}$
and $\mathfrak h=L_{N\otimes V_\kappa}$.
The symmetric convolution envelope is $U_\star\mathfrak g$.
Its PBW component with $a$ output blocks is
$\operatorname{Sym}^a_\star\mathfrak g$.

The relative deformation complex at $\alpha_0$ is
\begin{equation}\label{cbc:va:relative-convolution}
 \mathcal T=
 R\operatorname{Hom}
 \left(
 \operatorname{CE}_\star(\mathfrak h)/
             \operatorname{CE}_\star(\mathfrak g),
 U_\star\mathfrak g
 \right),
 \qquad d_{\mathcal T}=d+[\tau_0,-].
\end{equation}
Here $\operatorname{CE}_\star$ is the Chevalley
coalgebra on suspended inputs, and $\tau_0$ is the
twisting cochain of $\epsilon\otimes\mathrm{id}$.
Resolve coefficient Homs before taking this convolution
and the limit over configurations. Take products over
$G>0$. Convolution adds $G$.
The differential preserves $G$ and the PBW output number.

The bar--cobar comparison and
Lemma~\ref{cbc:va:chevalley-bar} identify
this complex with the relative deformation complex
of~\eqref{cbc:va:pointed-action}. In the Chevalley
comparison, the leading PBW filtration is the ordinary
symmetric--exterior Koszul complex. Its positive
factor count is bounded at each configuration, as in
Lemma~\ref{lem:words-arity-bound}. The comparison
therefore has no infinite length sum there.
The positive $G$ filtration makes the relative
convolution pronilpotent.

\begin{lemma}\label{cbc:va:action-degree}
In total $t$-degree $G$, the component of $\mathcal T$
with $a$ output blocks has a model in degrees at least
$a-2G$. For $G=1,a=2$, a degree-zero cocycle is an
ordinary binary coefficient, and it has no degree-zero
boundary.
\end{lemma}
\begin{proof}
With $n$ suspended inputs the source degree is
$2G-n$. A coefficient extension of degree $e$ into
a degree-zero target contributes degree $e+n-2G$.
Derived Hom between coefficient modules has a model
in extension degrees $e\geq0$.

Support forces $n\geq a$: restriction of an input
supported on at most $n$ blocks to the locus of $a$
distinct output blocks is zero when $n<a$.
Adjunction gives the same vanishing for derived Hom.
These statements give the lower bound.
For $G=1,a=2$, degree zero forces $n=2,e=0$.
The internal cocycle equation gives ordinary Hom,
and degree minus one is absent.
Products and the configuration limit preserve the
lower bound.
\end{proof}

Let $\ell:N\to k$ be the augmentation derivation
given by $\ell(t_n)=\delta_{n0}$ and zero on all
other polynomial degrees. It satisfies
\[
 \ell(Tb)=0,\qquad
 \ell(ab)=\ell(a)\epsilon(b)+\epsilon(a)\ell(b).
\]

On singleton suspended inputs define
\[
 \kappa_1\bigl(s(t_n\otimes v)\bigr)=\delta_{n0}v,
\]
and let it vanish elsewhere.
This is translation compatible: applying total
translation on the source gives the translation of
the displayed output. Its degree is minus one.
The augmentation-derivation identity makes it a
cocycle in $\mathcal T$.
Equivalently, the infinitesimal vertex map
$t_0\mapsto u$, $t_n\mapsto0$ for $n>0$, over
$|u|=2$, $u^2=0$, has twisting cochain
$\tau_0+u\kappa_1$.

There are no cochains of degree below minus one at
$G=1,a=1$, by Lemma~\ref{cbc:va:action-degree}.
Since $\kappa_1(s(t_0\otimes\mathbf1))=\mathbf1$,
its class is nonzero.

On the Heisenberg Fock space $k[x_1,x_2,\ldots]$,
let $J_1=\kappa\partial_{x_1}$.
The current mode identity and $J_0=0$ give
\[
 [J_1,Y(v,z)]=Y(J_1v,z),\qquad [T,J_1]=-J_0=0.
\]
Thus $J_1$ is a translation-compatible vertex
derivation, kills the vacuum and satisfies
$J_1J=\kappa\mathbf1$.

At two separated singleton blocks define
\begin{equation}\label{cbc:va:boundary-loop}
 h\bigl(s(t_n\otimes v),s(\mathbf1\otimes w)\bigr)
     =\delta_{n0}\,v\boxtimes J_1w.
\end{equation}
Extend by suspended symmetry, with all other
components zero. This has degree zero and is relative
to the boundary copy of $C$.

If $\widehat J_1=J_1\tau_0$, then
$h=-(\kappa_1*\widehat J_1)^{(2)}$ in convolution.
The minus sign cancels the sign from a degree-one
map crossing the suspended $t$ input of degree one.
Both factors are closed and the differential obeys
the Leibniz rule. Hence $h$ is closed.

\subsection{Extracting the separated vacuum coefficient}

For a component $K$ with $a$ output blocks, add a
separated boundary singleton at $z$ carrying $J$.
Apply the vacuum coefficient at that block and take
the constant coefficient of its Laurent expansion
at $z=\infty$. Use test function $1$ in $z$.
Put the suspended current first in the input order,
including that permutation's Koszul sign.
For a cochain of degree $q$, multiply the result by
$(-1)^{q+1}$ and denote it by $\mathsf P(K)$.

This extraction is defined on a coefficient resolution.
The input $J$ is a free translation generator: the
positive-weight Heisenberg states are free over $k[T]$,
and no translate has weight one. To see freeness,
translation is injective on nonconstant differential
polynomials by the largest-mode argument. Homogeneous
complements to its image then give a free $k[T]$-basis.
The same argument applies to $N_+$ and $N_+\otimes V_\kappa$.
On the latter, first use the highest actor weight.

\begin{lemma}\label{cbc:va:vacuum-probe}
The extraction has degree minus one and satisfies
\[
 d_{\mathcal T}\mathsf P+\mathsf P d_{\mathcal T}=0,
 \qquad \mathsf P(h)=\kappa\,\kappa_1.
\]
\end{lemma}
\begin{proof}
Extraction at the separated block commutes with
coefficient differentials, residues and derivatives
in every remaining position. In the component of a
fixed output partition, the differential acts within
one block. A bracket in the added singleton would
require a second input there. Removing its sole
input gives the zero component of that output block.
The two convolution terms multiplying by the added
$J$ cancel, because different output blocks commute
on their disjoint locus.

The remaining differential is the one on the other
blocks. Inserting the suspended current has degree
minus one. The stated factor $(-1)^{q+1}$ gives the
displayed anticommutation in the Hom differential.
Extraction preserves $G$ and finite normal order,
so it passes to their completed convolution.
On~\eqref{cbc:va:boundary-loop} its value is the
vacuum coefficient of $J_1J=\kappa\mathbf1$.
The permutation sign cancels its prefactor, giving
$\kappa\kappa_1$.
\end{proof}

Consider the pointed comparison induced by
\eqref{cbc:va:cochain-roof}:
\[
 \Theta_F:\operatorname{Map}_{J/}
                 (F,C_{\mathcal Z_\kappa})
                \longrightarrow\Phi_C(F,u).
\]
At $g_0$, its derivative sends source cochains to
$\mathcal T$. Every class in its $G=1$, two-output,
degree-zero image has zero vacuum probe.

Indeed, transfer the roof through the quartet
contraction. Its twisting corestriction is Lie-valued.
One can construct it by the finite rooted-tree
formulas for the bracket and the contraction.
Each corestriction has one output block.
In the derivative of composition, a source cochain
acts on one block and the base point acts on the others.

At a separated boundary singleton carrying $J$,
the resulting state is either $J$ or the contraction
of $\zeta_{(-1)}(J+p)$. Here $\zeta$ lies in the
resolving coefficient.
All states have nonnegative conformal weight.
The normal product adds weights, and the contraction
preserves them. Both possible values have positive
weight and hence zero vacuum coefficient.

This factorization through positive conformal weights
persists on coefficient resolutions.
Lemma~\ref{cbc:va:action-degree} places the indicated
degree-zero component in ordinary binary Hom, so
there is no hidden lower-degree correction to this
vanishing. It holds before extraction at infinity.
But $\mathsf P(h)=\kappa\kappa_1$ represents a
nonzero class. Thus $[h]$ is outside the image of
the derivative of $\Theta_F$.

\subsection{The coherent loop and its consequence}

Over polynomial forms $k[s,ds]$, the twisting cochain
$\tau_0+ds\,h$ is a loop at $\tau_0$.
It also defines a loop of coherent monoid actions.
For $p$ actor factors, set
\[
 \epsilon_p=\epsilon^{\otimes p},\qquad
 \ell_p=\sum_{i=1}^p
     \epsilon^{\otimes(i-1)}\otimes\ell
                    \otimes\epsilon^{\otimes(p-i)}.
\]
Let $\kappa_p$ be its singleton cochain and put
$h_p=-(\kappa_p*(J_1\tau_0^{(p)}))^{(2)}$.

Each is closed. Adjacent actor multiplication
preserves $\ell_p$ by its derivation identity,
and actor-unit insertion preserves it because
$\ell(1)=0$.
For composition of two actor words, the two linear
terms agree with
\[
 \ell_{p+q}=\ell_p\otimes\epsilon_q+
                         \epsilon_p\otimes\ell_q.
\]
Every term containing two copies of $ds$ is zero.

Use these formulas on the free associative bar
resolution of an actor word followed by its module
input. Each parenthesization of a word of length
$p$ gives the same constant term and the same
linear term $h_p$. Thus every face and degeneracy
agrees with its prescribed multiplication, action
or unit map. Constant fillers on these bar
simplices give all higher coherences.
The extra unit degeneracy contracts the augmented
bar resolution. The formulas therefore descend
to the required coherent action, retaining their
collision maps by Lemma~\ref{cbc:va:free-field-tensor}.

\begin{theorem}\label{cbc:va:cochain-nonfinal}
The comparison $\Theta_F$ is not surjective on
fundamental groups at $g_0$.
For any coherent transverse monoid structure on
$C_{\mathcal Z_\kappa}$ with action
$\alpha_{\mathcal Z}$, the mapping space from
$(C_N,\alpha_0)$ to that action is disconnected.
In particular this cochain action is not a local center.
\end{theorem}
\begin{proof}
For a nilpotent relative convolution Lie algebra,
based loops of its Maurer--Cartan space are
$\exp H^0$ with the Campbell--Hausdorff law.
The class of a constant loop form $ds\,v$ is
$\exp[v]$. See \cite{Getzler2009}*{\S5}.
Here one applies this statement at each finite
$G$ quotient. The differential preserves $G$,
so the quotients on cohomology are surjective.
Taking their inverse limit gives the completed
Campbell--Hausdorff group.

In composition of loops, terms involving two
deformation forms vanish. Hence $\Theta_F$ sends
$\exp[v]$ to the exponential of its derivative
on $[v]$. The nonzero probe of $[h]$ and the
vanishing probe on the image show that
$\exp[h]$ is not in that image.

The constructed loop belongs to
$\operatorname{Act}_{C_N}(C)$, with all its monoid
coherences. If it lifted to a loop of monoid maps
into $C_{\mathcal Z_\kappa}$, forgetting the
monoid structures would contradict the preceding
pointed comparison. The homotopy exact sequence
of~\eqref{cbc:va:action-fibre} therefore gives at
least two connected components in that fiber.
A final action would make it contractible.
\end{proof}

The probe uses a derivation that sends a positive
weight state to the vacuum. It detects this
specified action and its weight-preserving
resolutions. It does not exclude a final transverse
monoid with different coefficients or a
different action. That question requires recipients
for arbitrary monoid actions with their prescribed paths.

\section{Derivation profiles}
\label{cbc:va:profiles}

A derivation supplies an infinitesimal action. Extending
it to an action requires products of its values and
derivatives of its actor inputs. We construct those
terms together. Their order matters even when the
actor inputs themselves are symmetrized.

Let $\mathfrak g$ be a positive Lie object in the
symmetric coefficient category above, concentrated in
degree zero, and put $A=U_\star\mathfrak g$.
Let $D:\mathfrak g\to\mathfrak g$ be a degree-zero Lie
derivation that commutes with the full coefficient
maps. It extends to an associative derivation of $A$.
In particular, it commutes with collision and normal
transfer. A derivation on one separated fiber alone
does not meet this hypothesis.

Give $\varepsilon\mathfrak g$ degree one and form the
semidirect product with brackets
\[
 [x,\varepsilon m]=\varepsilon[x,m],
 \qquad [\varepsilon m,\varepsilon n]=0.
\]
Its suspended actor inputs $s\varepsilon m$ have
degree zero. Their number is the \emph{actor weight}.
The Chevalley coalgebra has underlying comodule
\[
 \operatorname{CE}_\star(\mathfrak g)\star
 \operatorname{Sym}^c_\star(\varepsilon\mathfrak g[1]),
\]
where the superscript $c$ specifies the unshuffle
coproduct. We identify the second factor's generators
with $m\in\mathfrak g$ in degree zero.

Write $P_w=\operatorname{Sym}^w_\star\mathfrak g$.
For $u=m_1\circ\cdots\circ m_w$, set
\[
 [x,u]=\sum_i
 m_1\circ\cdots\circ[x,m_i]\circ\cdots\circ m_w.
\]
Symmetric powers are coinvariants, identified with
invariants by averaging.

\begin{lemma}\label{cbc:va:profile-equation}
Coefficient maps $\theta_w:\mathfrak g\star P_w\to A$,
$w\geq1$, are the additional components of a twisting
cochain under $\operatorname{CE}_\star(\mathfrak g)$
if and only if
\begin{align}
0={}&[x,\theta_w(y;u)]-\theta_w(y;[x,u])
      -[y,\theta_w(x;u)]+\theta_w(x;[y,u])
      -\theta_w([x,y];u)\notag\\
 &+\sum_{\substack{S\sqcup T=[w]\\S,T\ne\varnothing}}
 \bigl(\theta_{|S|}(x;m_S)\theta_{|T|}(y;m_T)
       -\theta_{|S|}(y;m_S)\theta_{|T|}(x;m_T)\bigr).
 \label{cbc:va:profile-mc}
\end{align}
The splitting $(S,T)$ is ordered, and $m_S$ denotes
the symmetric product of the inputs with labels in $S$.
\end{lemma}
\begin{proof}
The value of a degree-one cochain on $n$ suspended
ordinary Lie inputs has degree $1-n$ in $A$.
Only $n=1$ contributes, and its value is
$\theta_w(x;u)$. The twisting equation has degree
$2-n$ and therefore only tests $n=2$.

Use $Q_2(sx\,sy)=s[x,y]$. Its three bracket terms are
$\theta_w([x,y];u)+\theta_w(y;[x,u])
-\theta_w(x;[y,u])$.

In a nonzero convolution product, each tensor factor
has one ordinary suspended input. Passing the second
degree-one cochain across the first factor contributes
a minus sign. The unshuffles with an empty actor part
give the two commutators, and those with both parts
nonempty give the sum in~\eqref{cbc:va:profile-mc}.
Multiplying the twisting equation by minus one proves
the formula. Weight zero is the envelope relation.
\end{proof}

\subsection{The fixed-point construction}

We first prove a universal identity in an associative
algebra with derivation. Let $z$ be a central formal
parameter, and for $R\in zA[[z]]$ set
\[
 f\langle R\rangle
   =\sum_{j\geq0}\frac{D^jf}{j!}R^j.
\]
Every power of $R$ stands on the right.
Choose central scalar parameters $t_1,t_2,\ldots$ and put
\begin{equation}\label{cbc:va:profile-fixed-point}
 U=\sum_{r\geq1}\frac{t_rz^r}{r!}m^r,\qquad
 R=U\langle R\rangle,\qquad G_x=x\langle R\rangle.
\end{equation}
Solve for $R$ by increasing powers of $z$.
Each unknown coefficient on the right is multiplied
by a coefficient of positive $z$-order, so the
solution is unique.

Let $\delta_x$ vary the actor variable $m$ in the
direction $[x,m]$. It replaces $D^jm$ by $D^j[x,m]$
and fixes all external arguments.

\begin{lemma}\label{cbc:va:profile-transport}
The series in~\eqref{cbc:va:profile-fixed-point} satisfy
\begin{align}
 \delta_xR&=[G_x,R],\label{cbc:va:profile-transport-r}\\
 [G_x,G_y]-G_{[x,y]}-\delta_xG_y+\delta_yG_x&=0.
 \label{cbc:va:profile-transport-g}
\end{align}
\end{lemma}
\begin{proof}
Put $U_j=D^jU/j!$ and $X_i=D^ix/i!$.
Since $\delta_xU=[x,U]$, differentiation of the
fixed-point equation gives a linear equation for
$Z=\delta_xR$:
\[
 Z=[x,U]\langle R\rangle+
       \sum_jU_j\sum_{p+q=j-1}R^pZR^q.
\]
Its linear part raises $z$-order, so its solution is
unique.

Taking the commutator of $G_x$ with the
fixed-point equation gives the same equation for
$[G_x,R]$. To check the constant term in that
equation, expand Leibniz for $D$. The difference is
\[
 \sum_{i,j}X_i[R^i,U_j]R^j
   =\sum_iX_i\sum_{p+q=i-1}R^pWR^q,
 \qquad
 W=\sum_j[R,U_j]R^j=[R,R]=0.
\]
This proves~\eqref{cbc:va:profile-transport-r}.

Write $Y_j=D^jy/j!$. Subtracting
$G_{[x,y]}=\sum_{i,j}[X_i,Y_j]R^{i+j}$ from
$[G_x,G_y]$ leaves
\[
 \sum_{i,j}X_i[R^i,Y_j]R^j
       -\sum_{i,j}Y_j[R^j,X_i]R^i.
\]
Expand each commutator with a power of $R$.
The first sum is $-\delta_yG_x$ and the second,
with its displayed minus sign, is $\delta_xG_y$,
by~\eqref{cbc:va:profile-transport-r}.
This gives~\eqref{cbc:va:profile-transport-g}.
\end{proof}

For a nonempty set $B$ of actor labels, define
\[
 u_B=\frac1{|B|!}\sum_{\sigma\in S_B}
                 m_{\sigma(1)}\cdots m_{\sigma(|B|)}.
\]
Let $\mathcal T_j$ be the nonnegative tuples
$(a,b_1,\ldots,b_j)$ such that
\[
 a+\sum_{i=1}^jb_i=j,\qquad
 a+\sum_{i=1}^\ell b_i\geq\ell+1
                    \quad(0\leq\ell<j).
\]
These are the child numbers of a plane rooted tree
with $j$ nonroot vertices, in preorder.
For arbitrary envelope elements $v_i$, put
\[
 T_j^D(x;v_1,\ldots,v_j)=
 \sum_{(a,b)\in\mathcal T_j}\frac{D^ax}{a!}
 \sum_{\sigma\in S_j}
 \frac{D^{b_1}v_{\sigma(1)}\cdots D^{b_j}v_{\sigma(j)}}
      {b_1!\cdots b_j!},
 \qquad T_0^D(x)=x.
\]

\begin{theorem}\label{cbc:va:profiles-theorem}
For every scalar sequence $\mathbf t=(t_1,t_2,\ldots)$,
the operations
\begin{equation}\label{cbc:va:profile-partitions}
 \theta_w^{\mathbf t}(x;m_1,\ldots,m_w)
 =\sum_{\pi\in\operatorname{Part}([w])}
     \left(\prod_{B\in\pi}t_{|B|}\right)
     T_{|\pi|}^D(x;u_B:B\in\pi)
\end{equation}
satisfy~\eqref{cbc:va:profile-mc}. The partitions here
are unordered and have nonempty blocks.
They define a map of derived convolution coalgebras
\[
 \operatorname{CE}_\star
          (\mathfrak g\ltimes\varepsilon\mathfrak g)
 \longrightarrow\operatorname{CE}_\star(\mathfrak g)
\]
under $\operatorname{CE}_\star(\mathfrak g)$.
On a configuration $I$, its nonconstant operations
use only $t_r$ with $r<|I|$.
\end{theorem}
\begin{proof}
Expanding the fixed point orders the subtrees below
each vertex. A vertex with $j$ children contributes
$D^jU/j!$, and the root contributes $D^ax/a!$.
Polarizing a term of $U$ groups the actor labels into
the blocks $B$, with product $u_B$ and coefficient
$t_{|B|}$. Ordering these blocks among the tree's
vertices and summing their orders gives $T_{|\pi|}^D$.

Consequently
\[
 G_x=x+\sum_{w\geq1}\frac{z^w}{w!}
              \theta_w^{\mathbf t}(x;m,\ldots,m).
\]

The coefficient of $z^w/w!$ in
\eqref{cbc:va:profile-transport-g} is
\eqref{cbc:va:profile-mc} on equal inputs.
The two variations give its actor substitutions.
Multiplication of exponential generating series
gives the ordered unshuffles. Polarization proves it for
all inputs.

The twisting cochain gives the map to $B_\star A$
by its iterated reduced coproduct. The
comparison of Lemma~\ref{cbc:va:chevalley-bar} identifies this target
with $\operatorname{CE}_\star(\mathfrak g)$.
On weight zero this is the canonical comparison,
which gives the specified map under that coalgebra.
Each positive input occupies a nonempty subset of
$I$. A term of actor weight $w$ also consumes one
ordinary input, so $w+1\leq|I|$. All coproduct and
tree sums are therefore finite on each configuration.
\end{proof}

The first two operations display both sources of
higher terms:
\begin{align*}
 \theta_1^{\mathbf t}(x;m)&=t_1Dx\,m,\\
 \theta_2^{\mathbf t}(x;m,n)
   &=t_2Dx\,\frac{mn+nm}{2}\\
   &\quad+t_1^2\left(
       D^2x\,\frac{mn+nm}{2}+Dx(Dm\,n+Dn\,m)\right).
\end{align*}
At weight three, the mixed coefficient of $t_1t_2$
is $\sum_{\{i\}\sqcup B=[3]}T_2^D(x;m_i,u_B)$.
Thus independently chosen leading weights also have
specified mixed primitives.

To interpret these maps as actions of a field
coalgebra, one needs a coefficient comparison.
Precisely, let $E=\Lambda(\varepsilon)$, with
$|\varepsilon|=1$ and $d\varepsilon=T\varepsilon=0$.
Suppose the residue construction identifies
$C_{E\otimes V}$ and $C_V$ with the two Chevalley
coalgebras in Theorem~\ref{cbc:va:profiles-theorem},
respectively. Require these equivalences to preserve
the vacuum axis, actor weights, collisions and normal
transfer, and to identify their derived mapping
spaces. Together with
Lemma~\ref{cbc:va:free-field-tensor}, they turn the
constructed map into a pointed action of $C_E$ on
$C_V$. Under these hypotheses, no completion of
normal derivatives or of state spaces is needed:
the formal parameter has proved finite identities
at every configuration.

\begin{lemma}\label{cbc:va:heisenberg-admissibility}
For $V=V_\kappa$, $\kappa\ne0$, these comparisons
exist. They are natural under rescaling
$\varepsilon\mapsto t\varepsilon$, tensor powers
of $E$, and finite scalar extension.
\end{lemma}
\begin{proof}
Resolve $E$ by the regular commutative field
algebra
\[
 \mathcal K=k[u,Tu,\ldots]\otimes\Lambda[c,Tc,\ldots],
 \qquad |u|=0,\quad |c|=1,\quad du=Tc,\quad dc=0.
\]
The map $q:\mathcal K\to E$ sends $u$ to zero,
$c$ to $\varepsilon$, and their positive
translates to zero.
The pairs $(T^nu,T^{n+1}c)$ contract by their
Euler homotopy, leaving only $\Lambda(c)$.
Thus $q$ is a quasi-isomorphism.

Apply the lattice comparison of
Lemma~\ref{cbc:va:free-field-tensor} to
$\mathcal K$ and $V_\kappa$.
The finite partition filtration preserves the
quasi-isomorphism induced by $q$, so it gives
$C_E\otimes^!C_{V_\kappa}\simeq C_{E\otimes V_\kappa}$.
The diagonal Lie object of $E\otimes V_\kappa$
is $\mathfrak g\ltimes\varepsilon\mathfrak g$:
the exterior coefficient multiplies the
chiral bracket, and $\varepsilon^2=0$.
This identifies the target with its Chevalley
model, including the vacuum axis.

The maps $u\mapsto tu$, $c\mapsto tc$ commute
with $q$ and with the lattice construction.
They give weight $t^w$ on $w$ suspended actor
inputs. Tensor powers use the same independent
contractible pairs. The constructions preserve
finite scalar extension and all collision maps,
since these maps already act on the lattices
and their residue resolutions.
\end{proof}

Henceforth in the Heisenberg case take
$D=\kappa^{-1}J_1=\partial_{x_1}$, so $D(J)=\mathbf1$.
It is the derivative of the vertex automorphism
$J\mapsto J+t\mathbf1$ and is therefore a map
of the full coefficient family.
Let $\alpha_{\mathbf t}$ denote the pointed
actions just constructed. Write $\alpha_D$ for
the sequence $t_1=1$, $t_r=0$ for $r>1$.

\begin{proposition}\label{cbc:va:leading-profiles}
For every $r\geq1$, the first variation of
$\alpha_{\mathbf t}$ in $t_r$ at the trivial
action is a nonzero class. The classes for
different $r$ have distinct actor weights.
For every $N\geq1$ there is a nontrivial pointed
action equal to the trivial action on all
configurations with at most $N$ positions.
\end{proposition}
\begin{proof}
The first variation is
$u_r(x;m_1,\ldots,m_r)=Dx\,u_{[r]}$.
In its PBW component with $r+1$ output roots,
it is $\gamma_r(x;v)=(Dx)\circ v$, where
$v\in P_r$ and $\circ$ is symmetric multiplication.
Its Chevalley cocycle equation is
\[
 [x,Dy]+[Dx,y]-D[x,y]=0.
\]
Thus it defines the triangular module extension
$P_{r+1}\oplus P_r$ with off-diagonal action
$\gamma_r$.

There is no coefficient map $P_r\to P_{r+1}$.
Indeed $P_{r+1}$ is open direct image from its
$(r+1)$-root locus, where the source restricts
to zero. The underlying extension therefore
has only its displayed splitting. It is not
equivariant: on $r+1$ distinct roots, put $J$
at the acting root and the vacuum at the others.
Then $\gamma_r$ gives the nonzero vacuum tensor
with coefficient one.

Degree-zero derived module maps here are ordinary
equivariant maps: their resolved coefficient Hom
and Chevalley complexes have no negative degrees.
The extension class is therefore nonzero.
PBW projection preserves the adjoint differential,
so $[u_r]$ is nonzero as well.

For the last assertion, use $t_N=1$ and all
other parameters zero. Its positive terms
require at least $N+1$ supports.

It remains to exclude a homotopy to the trivial
action. More generally, if two normalized
relative actions vanish below weight $r$,
write a gauge between them as $G=1+\sum_{j>0}g_j$.
The equation $d_0G=G\theta-\theta'G$ gives
$d_0g_j=0$ below weight $r$ and
$d_0g_r=\theta_r-\theta'_r$ in weight $r$.
Their first classes must agree. Here one of
them is $[u_N]\ne0$ and the other is zero.
\end{proof}

This also obstructs the ordinary exterior product
for profiles starting above weight one.
Let $\beta$ have first nonzero class
$[\beta_r]$ in weight $r>1$.
If the standard product $\mu$ of $C_E$ made it
multiplicative, pull its action path back along
the primitive map $\varepsilon\mapsto
\varepsilon_1+\varepsilon_2$.
On the product side this scales the first
class by $2^r$. On the iterated-action side
the first jets add, giving $2[\beta_r]$:
all products of positive lower jets vanish.
The preceding gauge calculation would give
$(2^r-2)[\beta_r]=0$, impossible in
characteristic zero.

This concerns the standard multiplication
on this actor. It does not exclude a map of
the action to a different multiplicative action.

\begin{example}\label{cbc:va:quadratic-profile}
In $k[q]$, take $D=\partial_q$, $x=m=q$ and
$t_2=1$, with all other parameters zero.
Then $G_q=q+R$ and
\[
 G_q=q+\frac{z^2}{2}G_q^2.
\]
The solution with constant term $q$ is
\[
 G_q=\sum_{n\geq0}\frac1{n+1}\binom{2n}{n}
                    \left(\frac{z^2}{2}\right)^nq^{n+1}.
\]
Indeed, $G_q=qC(z^2q/2)$ reduces the equation to
$C=1+uC^2$, whose solution at $C(0)=1$ is
$(1-\sqrt{1-4u})/(2u)$.
\end{example}

\subsection{An inner gauge and its domain}
\label{cbc:va:inner}

The same fixed point determines when an inner
derivation gives a trivial profile. We first work
in an ordinary unital associative algebra $A$.
Fix $a\in A$ and take $D=[a,-]$.

The variations $\delta_x$ fix $a$.
They satisfy $[\delta_x,\delta_y]=-\delta_{[x,y]}$,
so $\rho_xF=[x,F]-\delta_xF$ defines a Lie action
on functions of $m$. Let $d_0$ be its Chevalley
differential. In particular, $d_0F(x)=\rho_xF$ on
zero-cochains.

The positive one-cochain $\theta(x)=G_x-x$ satisfies
$d_0\theta+\theta\smile\theta=0$ by
\eqref{cbc:va:profile-transport-g}.
A normalized gauge from $\theta$ to zero is an
invertible zero-cochain $B=1\bmod z$ with
$d_0B=B\theta$.

\begin{proposition}\label{cbc:va:inner-gauge}
For the fixed point~\eqref{cbc:va:profile-fixed-point}
with $D=[a,-]$, such a gauge is
\begin{equation}\label{cbc:va:ordered-gauge}
 B=\sum_{n\geq0}\frac{(-1)^n}{n!}a^nR^n.
\end{equation}
Its powers of $a$ stand to the left of its powers
of $R$.
\end{proposition}
\begin{proof}
The constant term is one, so its geometric inverse
converges $z$-adically.
Since $\delta_xa=0$, transport gives
\[
 \delta_xB
   =\sum_{n\geq0}\frac{(-a)^nG_xR^n}{n!}-BG_x.
\]
The coefficient on the left of $R^k$ in the sum is
\[
 \sum_{n+j=k}\frac{(-a)^n(\operatorname{ad}a)^jx}{n!j!}
       =\frac{(-1)^k}{k!}xa^k.
\]
For this equality, left and right multiplication
$L_a,R_a$ commute and $\operatorname{ad}a=L_a-R_a$.
The coefficient identity is exactly
$e^{-tL_a}e^{t(L_a-R_a)}=e^{-tR_a}$.
Thus $\delta_xB=xB-BG_x$, and
$[x,B]-\delta_xB=B(G_x-x)$, as required.
\end{proof}

For $U=zm$ the first terms are
\[
 R=zm+z^2[a,m]m+O(z^3),\qquad
 B=1-zam+z^2\left(-a[a,m]m+\frac12a^2m^2\right)
       +O(z^3).
\]
The coefficient of $z^2$ in $B-e^{-aR}$ is
$\frac12a[a,m]m$, nonzero in the free algebra on
$a,m$. An ordinary exponential does not give this
gauge.

The polarized identities also hold in a resolved
coefficient convolution under two hypotheses.
All operations must preserve that coefficient family,
and sufficiently many positive factors must vanish
at each configuration.
The inverse and the gauge equation are then finite
there. These hypotheses include the element that
implements the inner derivation.

For the positive convolution envelope
$U_\star\mathfrak g$, that last hypothesis has only
the zero-derivation instance. Its convolution unit
$K$ is supported at the empty configuration, whereas
$\mathfrak g$ is positive. A global element
$a:K\to U_\star\mathfrak g$ is therefore scalar,
and $[a,-]=0$. A nonzero state mode can still define
a derivation of the full coefficient family. It
need not arise from a global element of this
envelope. In particular, the gauge above does not
trivialize the Heisenberg shift used in the vacuum
probe.

\section{Multiplicative recipients for the profiles}
\label{cbc:va:profile-recipients}

The failure of the standard product on a higher
profile does not prevent that action from mapping
to a multiplicative one. We construct such maps
by changing the actor parameter. To compare their
actions, the change of parameter must extend to
an actual bar map.

\subsection{Tensoring the entire coefficient construction}

\begin{lemma}\label{cbc:va:strong-tensor}
Let $C$ be a cocommutative factorizing
$\star$-coalgebra with empty coefficient $k$.
Then $\mathcal T_C(F)=F\otimes^!C$ is a
colimit-preserving strong symmetric monoidal
functor for $\star$.
For $C=C_V$, it takes $L_k$ to $L_V$ and
$L_E$ to $L_{E\otimes V}$, with their brackets.
\end{lemma}
\begin{proof}
Use the coproduct of $C$ to compare
$(F\star G)\otimes^!C$ with
$(F\otimes^!C)\star(G\otimes^!C)$.
On the summand for $I=A\amalg B$, this is
open direct image of
\[
 (F_A\boxtimes G_B)\otimes^!C_I|_{U_{A,B}}
 \longrightarrow
 ((F_A\otimes^!C_A)\boxtimes
                  (G_B\otimes^!C_B))|_{U_{A,B}}.
\]
It is an equivalence by factorization and the
open projection formula. That formula follows
here from localization: extraordinary restriction
to the complement vanishes, and restriction to
the open is the displayed tensor.
At the empty configuration the unit comparison
is $K\otimes^!C=K$.

Every iterated comparison uses the same final
partition and the same iterated coproduct.
Coassociativity and cocommutativity give all
monoidal coherences. Collision kills exactly
the summands whose blocks split a collision
fiber. On the others the comparison is the
same base-change map. Tensor product and
convolution preserve colimits.

On a small-diagonal branch, projection for its
full normal direct image identifies
$L_k\otimes^!C_V$ with the state coefficient
$L_V$. The bracket on $L_k$ is the vacuum
principal-part boundary. Tensoring it with
the factorization lattice of $V$ gives its
two-position residue boundary, hence the
bracket of $L_V$.
The same computation on the even and odd
summands of $L_E$ gives
$L_{E\otimes V}=\mathfrak g\ltimes\varepsilon\mathfrak g$.
Thus the functor transports the Chevalley,
envelope and bar constructions, including
their coefficient maps.
\end{proof}

For any degree-zero positive Lie object
$\mathfrak l$, work in
$U_\star(\mathfrak l\ltimes\varepsilon\mathfrak l)$
and set
\[
 \rho_r(m_1,\ldots,m_r)
   =\frac1{r!}\sum_{\sigma\in S_r}
      m_{\sigma(1)}\cdots m_{\sigma(r-1)}
                       (\varepsilon m_{\sigma(r)}).
\]

Define a cochain by $\tau_{\mathbf t}(sx)=x$
on ordinary singleton inputs and
\[
 \tau_{\mathbf t}
   (s\varepsilon m_1\cdots s\varepsilon m_r)
       =t_r\rho_r(m_1,\ldots,m_r)
\]
on pure actor inputs, with all other reduced
components zero.

It is a twisting cochain. With one ordinary
input its equation is adjoint equivariance.
With two ordinary inputs it is the envelope
relation. The remaining equation polarizes
from equal actor inputs, where
$\rho_r(m,\ldots,m)=m^{r-1}\varepsilon m$.
All products of two such values vanish because
$[m,\varepsilon m]=0$ and $(\varepsilon m)^2=0$.

Thus $\tau_{\mathbf t}$ gives the strict bar map
\begin{equation}\label{cbc:va:actor-bar-map}
 F_{\mathbf t}
   =\sum_{j\geq0}(s\tau_{\mathbf t})^{\star j}
                                  \overline\Delta^{(j)}.
\end{equation}
After the inverse of
Lemma~\ref{cbc:va:chevalley-bar}, it defines a
pointed self-map. For $\mathfrak l=L_k$ denote
it by $\Sigma_{\mathbf t}:C_E\to C_E$.
Augmentation kills its positive odd outputs,
so it also preserves the augmentation to $J$.
Every sum is finite on each configuration.

\subsection{A bar lift of the derivation action}

The substitution identity has a differential
version. Let $(\mathcal G,d)$ be a complete
filtered dg associative algebra, with
$F^p\mathcal G\,F^q\mathcal G\subseteq F^{p+q}\mathcal G$.
Let $D$ preserve this filtration and commute
with $d$. Suppose $a$ has degree one,
$da+a^2=0$, and $p\in F^1\mathcal G$ has
degree zero with $dp+[a,p]=0$.
Put
\[
 a_j=D^ja/j!,\qquad p_j=D^jp/j!,\qquad
 R=\sum_{j\geq0}p_jR^j,\qquad
 A_R=\sum_{j\geq0}a_jR^j.
\]

\begin{lemma}\label{cbc:va:covariant-substitution}
These series satisfy
$dR+[A_R,R]=0$ and $dA_R+A_R^2=0$.
\end{lemma}
\begin{proof}
The differentiated equations are
$da_j=-\sum_{u+v=j}a_ua_v$ and
$dp_j=-\sum_{u+v=j}[a_u,p_v]$.
Expanding the fixed point gives
\[
 \sum_{u,v}[a_u,p_v]R^{u+v}
 =[A_R,R]-\sum_vp_v[A_R,R^v].
\]
The difference between the two sides is
$\sum_u a_u[R^u,\sum_vp_vR^v]=0$.

Consequently $-[A_R,R]$ solves the linear
fixed-point equation obtained by differentiating
$R$. Its positive filtration gives uniqueness.
Since the $a_j$ are odd, substitution yields
\[
 dA_R=-\sum_{u,v}a_ua_vR^{u+v}
                         +\sum_u a_u[A_R,R^u].
\]
Subtracting the first sum from $A_R^2$ gives
$\sum_u a_u[R^u,A_R]$, which cancels the
last sum. This proves both equations.
\end{proof}

Apply the lemma with
$B=U_\star(\mathfrak g\ltimes\varepsilon\mathfrak g)$,
$A=U_\star\mathfrak g$ and the bar convolution
$\mathcal G=\operatorname{Hom}^\bullet(B_\star B,A)$.
Let $\pi:B\to A$ kill the odd ideal.
The degree-minus-one projection derivation
$q:B\to A$ is determined by
\[
 q(x)=0,\quad q(\varepsilon m)=m,\quad
 q(uv)=q(u)\pi(v)+(-1)^{|u|}\pi(u)q(v).
\]
It respects the envelope relations.
For example, applying $q$ to
$[x,\varepsilon m]=\varepsilon[x,m]$
gives the same $[x,m]$ on both sides.
Two-odd-input relations give zero.

Take the length-one cochains induced by $\pi$
and $q$ for $a$ and $p$. Their degrees are
one and zero, and the projection derivation
identity is $dp+[a,p]=0$.
The lemma gives a twisting cochain and a
coalgebra map
$\mathcal F_D:B_\star B\to B_\star A$.
On a word $[v_0|\cdots|v_w]$, its corestriction is
\begin{equation}\label{cbc:va:bar-tree}
 \sum_{(a,b)\in\mathcal T_w}
 \frac{D^a\pi(v_0)}{a!}
 \frac{D^{b_1}q(v_1)\cdots D^{b_w}q(v_w)}
      {b_1!\cdots b_w!}.
\end{equation}

This follows by the same ordered-subtree
expansion as the profile fixed point.
It vanishes unless $v_0$ has odd-generator
degree zero and each later $v_i$ has
odd-generator degree one.

Restriction along the Chevalley comparison
puts the ordinary input first and sums the
orders of the actor inputs, giving $\alpha_D$.
The bar-length filtration is finite at each
configuration, so this is an actual coefficient
map there.

\begin{proposition}\label{cbc:va:profile-comparison}
For the Heisenberg actions, there is a natural
pointed path
\[
 \alpha_D(\Sigma_{\mathbf t}\otimes^!\mathrm{id}_C)
                        \simeq\alpha_{\mathbf t}.
\]
It preserves all collisions, normal transfers
and the specified vacuum path.
\end{proposition}
\begin{proof}
Lemma~\ref{cbc:va:strong-tensor} transports
$\Sigma_{\mathbf t}$ from $L_k$ to $\mathfrak g$.
Represent the composite by
$\mathcal F_DF_{\mathbf t}$.
A nonzero term puts the ordinary input in
its first singleton slot and partitions the
actor inputs into $\rho$-blocks.
The projection derivation satisfies
$q\rho_r(m_1,\ldots,m_r)=u_{[r]}$.
For each unordered partition, summing its
block orders in~\eqref{cbc:va:bar-tree} gives
exactly~\eqref{cbc:va:profile-partitions}, with
coefficient $\prod_Bt_{|B|}$.

Thus the corestrictions agree in every weight.
The pointed bar comparisons transport this
equality to the stated path and its coherences.
All operations are maps of the full coefficient
families.
\end{proof}

\subsection{The complex of natural higher comparisons}

It remains to construct multiplication coherences
for $\alpha_D$ itself. The relevant primitives
can be chosen as universal differential
polynomials. No contraction of an unspecified
geometric mapping complex is assumed.

Let $P,Z$ be finite-dimensional degree-zero
spaces. Write $E(P)=\Lambda(P^*[-1])$ and
choose commuting derivations $D_1,\ldots,D_{\dim Z}$
of a degree-zero Lie algebra $\mathfrak g$.
A natural cochain on
$\operatorname{CE}(\mathfrak g\otimes E(P))$
is a universal multilinear differential polynomial
with values in $U(\mathfrak g)$.
It applies the chosen derivations and multiplies
the resulting inputs, using each input once. Its evaluation
includes suspended symmetrization.
Brackets are expanded as commutators.

Let $\mathcal N(P;Z)$ be their convolution
algebra, with the Hom differential twisted by
the scalar augmentation action $\tau$.
Take the product over actor weights, with
finitely many derivative orders at each
cochain degree and actor weight.
Each specified input is evaluated by a
finite polynomial.

For an exterior subset $S$ and derivative
multi-index $a$, let $x_{S,a}$ extract that
exterior coefficient and apply $D^a/a!$.
It has degree $1-|S|$, actor weight $|S|$
and derivative weight $|a|$.
These unary cochains freely generate
$\mathcal N(P;Z)$ under convolution.
Indeed their products are signed orbit
sums of ordered differential monomials.
Evaluation in the free differential Lie
algebra on distinct named inputs recovers
each ordered jet word in its free tensor
envelope, proving independence.
Polarization handles repeated exterior types.

Put $t=x_{\varnothing,0}=\tau$ and
$y_{S,a}=(-1)^{\binom{|S|}{2}}x_{S,a}$.
For disjoint subsets let $\epsilon(B,C)$
be the sign sorting their exterior product.
The differential is
\begin{equation}\label{cbc:va:natural-differential}
 \delta t=t^2,\qquad
 \delta y_{S,a}
 =-\sum_{\substack{B\sqcup C=S\\b+c=a}}'
    (-1)^{|B|}\epsilon(B,C)y_{B,b}y_{C,c}.
\end{equation}
The prime omits $(\varnothing,0)$ on either
side, and the second formula excludes $t$.
To obtain it, use
$Q_2(s(e_Bx)s(e_Cy))=(-1)^{|B|}
s(e_Be_C[x,y])$, the divided Leibniz rule,
and the convolution sign
$(-1)^{|g||c'|}$ in $(f*g)(c)$.
The terms containing $t$ cancel the
commutator with the scalar twisting cochain.

\begin{lemma}\label{cbc:va:natural-cohomology}
Weight by weight,
\[
 H^\bullet\mathcal N(P;Z)
   =\widehat{\operatorname{Sym}}(P)
                         \otimes\Lambda(Z[-1]).
\]
The $P$ generators have degree zero and actor
weight one. The $Z[-1]$ generators have
degree one and derivative weight one.
Thus actor weight $q$ and derivative order
$q$ have cohomology
$\operatorname{Sym}^q(P)\otimes\Lambda^q(Z)$
in degree $q$.
\end{lemma}
\begin{proof}
Let $\Gamma(Z)$ be the divided-power coalgebra,
with $\Delta z^{[a]}=\sum_{b+c=a}
z^{[b]}\otimes z^{[c]}$.
Formula~\eqref{cbc:va:natural-differential}
identifies $\mathcal N(P;Z)$ with
\[
 (k\langle t\rangle,\delta t=t^2)
        *\Omega\bigl(\Lambda(P[1])\otimes\Gamma(Z)\bigr).
\]
The augmentation ideal of the first factor
contracts by $t^{2j}\mapsto t^{2j-1}$ and
$t^{2j-1}\mapsto0$.
On a free-product word, contract its first
such factor with the sign of the prefix.
Projection killing $t$ is therefore a
quasi-isomorphism.

Map the remaining cobar to
$\operatorname{Sym}(P)\otimes\Lambda(Z[-1])$
by preserving its primitive generators and
killing the others. The differential becomes
a graded commutator or an odd square, hence
zero. The map is surjective on cohomology.

To prove injectivity, dualize each auxiliary
weight. The dual coalgebra algebra is
$B=\Lambda(P^*[-1])\otimes\operatorname{Sym}(Z^*)$.
With generators $\alpha_i,z_j$, its free
resolution of $k$ is
\[
 B\otimes k[v_i]\otimes\Lambda(\eta_j),
 \qquad |v_i|=0,\quad|\eta_j|=-1,\quad
 dv_i=\alpha_i,\quad d\eta_j=z_j.
\]

For an odd pair,
$d(v_i^{m+1})=(m+1)\alpha_i v_i^m$,
leaving only the scalar class.
An even pair is the two-term resolution
for multiplication by $z_j$.
Their tensor product resolves $k$.
Its Hom into $k$ has zero differential and
the stated symmetric--exterior dimensions.

The weightwise dual bar computes the same
Ext as the cobar above, so the surjection
has equal finite dimensions in every weight.
It is an isomorphism.
\end{proof}

For one derivation, every equal-weight sector
with $q\geq2$ is acyclic. For two derivations,
the degree-two sector at $q=2$ need not vanish.
The number of derivations matters to this
construction.

We also need the natural associative-bar
cochains. At each primitive PBW degree,
allow PBW projections followed by the same
ordered differential polynomials. Take their
product over primitive degrees.
This class is closed under composition and
extension of the exterior coefficients.
An actor face sets one exterior generator
to zero. Cochains vanishing on every actor
face have positive weight in every actor
coordinate.

Restriction along the Chevalley--bar map
is a surjective quasi-isomorphism of these
natural convolution complexes, also relative
to any union of actor faces.
Here is a construction of its kernel contraction.
For an abelian Lie algebra and fixed
multilinear primitive degree, the bar is
the finite complex of ordered nonempty
partitions, with adjacent-block merging.
The Koszul resolution in the preceding proof
identifies its homology with the suspended
symmetrization image. Choose rational
complements to cycles and boundaries and
average their projections over permutations.
This gives natural maps $p,h$ with
$pj=1$, $hj=ph=h^2=0$ and $dh+hd=1-jp$.

For a nonabelian Lie algebra, the difference
of the PBW product from the symmetric product
lowers primitive degree. If $b$ is the
resulting bar perturbation, set
$p'=p(1+bh)^{-1}$ and $h'=h(1+bh)^{-1}$.
The series are finite on every input.
The inclusion remains $j$, since $bj=jQ_2$
and $hj=0$. Its transferred differential
is precisely the Chevalley differential.

Precomposition with $p'$ splits restriction
on untwisted Hom complexes, and
$f\mapsto(-1)^{|f|}fh'$ contracts its kernel.
Twisting by the scalar action preserves
that kernel and raises minimum primitive
input degree. A second finite perturbation
therefore contracts the twisted kernel.
All maps preserve actor and derivative
weights and commute with the actor faces.

\begin{lemma}\label{cbc:va:coherent-fillings}
Fix one derivation and at least two actor coordinates.
Prescribe compatible natural action data on all proper
actor faces and the boundary of a parameter simplex or cube.
These data extend to the interior, and their extension
space is contractible.
The cochains here have derivative order
equal to actor weight.
\end{lemma}
\begin{proof}
Commuting actor-face projections give a
surjective restriction by inclusion-exclusion.
Its kernel consists of the cochains using
every actor coordinate.
Each of its weights is at least two,
so it is acyclic by
Lemma~\ref{cbc:va:natural-cohomology} and the
relative bar comparison just constructed.

Polynomial forms extend compatible boundary
values on a simplex or cube. Choose a
linear lift of the boundary and actor-face
data, and solve Maurer--Cartan successively
in actor weight. At the next weight its
curvature is a closed degree-two element
of the fully mixed kernel tensored with
relative polynomial forms. Closedness is
the Bianchi identity. The relative condition
is the already satisfied boundary equation.

This complex is acyclic over $k$, so a
degree-one correction removes the curvature
without changing the prescribed data.
The weight filtration is complete.

The same argument with extra parameter
simplices solves every relative lifting
problem. Thus the restriction of
Maurer--Cartan spaces is a trivial Kan
fibration and its fibers are contractible.
\end{proof}

\begin{theorem}\label{cbc:va:multiplicative-recipient}
The Heisenberg action $\alpha_D$ admits
a unital homotopy-coherent action of the
standard monoid $(C_E,\mu)$.
Every profile action $\alpha_{\mathbf t}$
maps to this multiplicative action by
$\Sigma_{\mathbf t}$.
Its two prescribed unit insertions have
the common factorizing recipient
\[
 \mu(\Sigma_{\mathbf t}\otimes^!\Sigma_{\mathbf t}):
 C_E\otimes^!C_E\longrightarrow C_E
\]
with target action $\alpha_D$.
\end{theorem}
\begin{proof}
Use the bar lift $\mathcal F_D$ for the
unary action. Its terms have derivative
order equal to actor weight.

For $n$ actors, a coherent monoid action
has a family on the cube $I^{n-1}$.
The face $t_i=0$ multiplies the adjacent
actors before applying the shorter family.
The face $t_i=1$ composes the two shorter
bar actions.
The unit at an end deletes an adjacent
parameter. At an interior position it
combines the two parameters by
$a\vee b=a+b-ab$.
This polynomial connection is associative,
has identity zero and absorbing element
one, so repeated unit deletions agree.

For $n=2$, these rules prescribe the
endpoints and the two unit faces.
They coincide where they meet.
Lemma~\ref{cbc:va:coherent-fillings} supplies
the binary path. Inductively, associativity
of actor multiplication and of bar-map
composition makes all prescribed boundary
intersections agree. The same lemma fills
each higher cube with these values fixed.
It also proves contractibility of the
natural fillings after the unary action
and a binary path are fixed.

Lemma~\ref{cbc:va:strong-tensor} identifies
extension by another exterior actor with
geometric tensoring of the action.
It carries $D$ to $D\otimes1$ by naturality
in $J\mapsto J+h\mathbf1$ over $k[h]/(h^2)$.
Thus the bar endpoints are the required
geometric endpoints, including their unit,
collision and normal-transfer data.
Every universal formula has finitely many
primitive inputs at a fixed configuration.

Proposition~\ref{cbc:va:profile-comparison}
supplies the profile maps.
Compose two of them with the constructed
multiplicative action. On either unit
axis, its unit path reduces to the
profile comparison. At the intersection
both restrict to the same pointing.
The binary coherence supplies the
relative two-simplex between these paths.
Higher coherences give the same assertion
for finite ordered products.
\end{proof}

Contractibility above concerns the specified
class of natural differential-polynomial
fillings. It is not a classification of
all geometric multiplicative structures
or of all pointed actions. The general monoid
recipient problem still requires their
full mapping spaces.

\section{Relations between separated inputs}
\label{cbc:va:separated}

Separation imposes two different equations on an
action. Operators from different blocks commute,
and a prescribed mixed input becomes their product.
Their primitives must in turn satisfy higher
equations. We construct a resolution that keeps
both kinds of relation.

Let $V_1,\ldots,V_r$ be vector spaces in degree zero.
All tensor products in this section are algebraic
over $k$. For a nonempty subset $S\subseteq[r]$, put
\[
 V_S=\bigotimes_{i\in S}^{\mathrm{increasing}}V_i,
 \qquad
 F_r=T\left(\bigoplus_{\varnothing\ne S\subseteq[r]}V_S\right),
 \qquad P_r=\bigotimes_{i=1}^rT(V_i).
\]
Here $T$ is the unital tensor algebra.

Write $x_i$ for its singleton generating maps and
$c_S$ for the other generating maps of $F_r$.
Every expression below is multilinear in its
indicated inputs. The quotient
$\vartheta_r:F_r\to P_r$ sends
\[
 c_S\longmapsto p_S=x_{i_1}\cdots x_{i_s},
 \qquad i_1<\cdots<i_s,
\]
and sends $x_i$ to the $i$th tensor factor.

For two one-dimensional spaces, the two relations
are already independent. Starting with
$k\langle x_1,x_2,c\rangle$, adjoin $h,e$ in degree
minus one and set
\[
 dh=x_1x_2-x_2x_1,\qquad de=c-x_1x_2.
\]
In the bases $(h,e)$ and $(x_1x_2,x_2x_1,c)$
the differential is
$\left(\begin{smallmatrix}1&-1\\-1&0\\0&1\end{smallmatrix}\right)$.
One primitive cannot impose both equations.

\subsection{The subset resolution}

Let $C_i=k\oplus V_i[1]$ be the coaugmented coalgebra
with primitive reduced part, and form
$C=\bigotimes_iC_i$. Its cobar algebra $R_r=\Omega C$
is free on maps
\[
 q_S:V_S\longrightarrow R_r^{1-|S|},
 \qquad q_{\{i\}}=x_i.
\]
For disjoint subsets define
$\operatorname{inv}(L,R)=\#\{(\ell,j)\in L\times R:\ell>j\}$.
Use the cobar convention
$d(tc)=-t(dc)+\sum(-1)^{|c'|+1}(tc')(tc'')$,
where $t$ has degree one.

\begin{theorem}\label{cbc:va:subset-resolution}
The differential of $R_r$ is
\begin{equation}\label{cbc:va:subset-differential}
 dq_S=\sum_{\substack{L\amalg R=S\\L,R\ne\varnothing}}
       (-1)^{|L|+1+\operatorname{inv}(L,R)}q_Lq_R.
\end{equation}
It squares to zero. The map
$\epsilon:R_r\to P_r$ that preserves the singleton
generators and kills $q_S$ for $|S|>1$ is a
quasi-isomorphism.
\end{theorem}
\begin{proof}
Each factor of $V_i[1]$ is odd. Unshuffling a subset
into $(L,R)$ contributes
$(-1)^{\operatorname{inv}(L,R)}$. Desuspension
contributes $(-1)^{|L|+1}$. These are the two signs
in~\eqref{cbc:va:subset-differential}.

In $d^2q_S$, each ordered three-part decomposition
of $S$ occurs twice, by splitting the first or the
second factor. The derivation sign makes their
coefficients opposite. Thus $d^2=0$.

Under $\epsilon$, the pair differential becomes a
commutator of different tensor factors. Every
larger split contains a higher generator, so
$\epsilon$ is a chain map.

Put $A_i=T(V_i)$. The sequence
\[
 0\longrightarrow A_i\otimes V_i
 \xrightarrow{a\otimes v\mapsto av}A_i
 \longrightarrow k\longrightarrow0
\]
resolves the augmentation module: every nonempty
word has a unique last factor. Tensoring these
resolutions gives $P_r\otimes_\tau C$, where
$\tau$ sends suspended singletons to their
generators and all larger subsets to zero.

Map $C$ to $BP_r$ by the signed shuffle of its
singleton inputs into bar slots. Interior faces
cancel in pairs because distinct tensor factors
commute. Splitting a shuffle at a bar cut proves
the coalgebra identity. Its $P_r$-linear extension
is a map of free resolutions of $k$ inducing the
identity on $k$, hence a chain homotopy equivalence.
Tensoring over $P_r$ with $k$ shows that
$C\to BP_r$ is a quasi-isomorphism.

Assign weight one to each singleton input.
At a fixed weight, reduced bar and cobar words have
bounded length. The cobar-length filtration now
shows that $\Omega C\to\Omega BP_r$ is a
quasi-isomorphism. The bar--cobar contraction of
$\Omega BP_r\to P_r$ is finite in that weight.
Their composite is $\epsilon$. Taking the direct
sum over weights proves the assertion.
\end{proof}

The first higher equation is the six-term relation
\begin{equation}\label{cbc:va:triple-differential}
 dq_{123}=x_1h_{23}-x_2h_{13}+x_3h_{12}
             -h_{12}x_3+h_{13}x_2-h_{23}x_1,
 \qquad h_{ij}=q_{\{i,j\}}.
\end{equation}
Write $k_{ijk}=q_{\{i,j,k\}}$. At four inputs,
\begin{align}
 dq_{1234}={}&x_1k_{234}-x_2k_{134}+x_3k_{124}-x_4k_{123}
                  \notag\\
 &+k_{123}x_4-k_{124}x_3+k_{134}x_2-k_{234}x_1
                  \notag\\
 &-h_{12}h_{34}+h_{13}h_{24}-h_{23}h_{14}
       -h_{14}h_{23}+h_{24}h_{13}-h_{34}h_{12}.
 \label{cbc:va:four-differential}
\end{align}
The pair--pair terms are necessary for closedness.

In multidegree $(1,\ldots,1)$, a word of generators
is an ordered partition of $[r]$. A partition into
$a$ blocks has degree $a-r$. These are the faces
of the permutohedron with their incidence
differential. For $r=3$ their dimensions are
$(1,6,6)$ in degrees $(-2,-1,0)$. For $r=4$
they are $(1,14,36,24)$.
The algebra also includes repeated inputs of each
color, retaining their order within that color.

There is a finite contraction in every multiweight.
Fix $\mathbf n=(n_1,\ldots,n_r)$.
After grouping coefficient factors by color, each
word with these counts has coefficient
$\bigotimes_iV_i^{\otimes n_i}$.
The differential is an integer matrix on a finite
set of words and preserves within-color order.

Choose the color-increasing word as the degree-zero
representative. In each degree, take in lexicographic
order the differential columns that increase rank.
Their source vectors complement the cycles. Their
images form a basis of boundaries.
Together with the chosen degree-zero representative,
these give bases of boundaries, representatives
and source lifts, by
Theorem~\ref{cbc:va:subset-resolution}.

Let $h$ send a boundary basis vector to its chosen
source lift and vanish on the other basis vectors.
If $i$ includes the color-increasing word, finite
matrix inversion gives
\[
 dh+hd=1-i\epsilon,\qquad
 h^2=hi=\epsilon h=0.
\]
The scalar matrices depend only on $\mathbf n$, so
this construction is natural in $V_i$.
It is a contraction of complexes, without a
multiplicativity assertion.

For complex inputs one includes their internal
differentials and the full Koszul signs in
$\Omega\bigotimes_i(k\oplus V_i[1])$.
The tensor-resolution argument still applies over
$k$. The scalar signs in
\eqref{cbc:va:subset-differential} are the
degree-zero specialization.

\subsection{Preserving the mixed generators}

The map $\vartheta_r$ also has prescribed values
on $c_S$. To retain them, adjoin to $F_r$ the
generators $q_S$ of degree $1-|S|$ and $e_S$ of
degree minus one for $|S|\geq2$. Call the resulting
free graded algebra $Q_r$ and put
\[
 dc_S=0,\qquad de_S=c_S-p_S,
\]
with $dq_S$ as in~\eqref{cbc:va:subset-differential}.

\begin{theorem}\label{cbc:va:relative-resolution}
The factorization $F_r\longrightarrow Q_r
\xrightarrow{\sim}P_r$ of $\vartheta_r$ is
relative semifree. The second arrow sends $c_S$
to $p_S$ and kills all negative-degree generators.
\end{theorem}
\begin{proof}
Adjoin the $q_S$ by increasing cardinality.
Their differentials use earlier generators. Each
$de_S$ lies in $F_r$. This proves relative
semifreeness and $d^2=0$.

Change generators from $c_S$ to $y_S=c_S-p_S$.
Then $de_S=y_S$ and $dy_S=0$, so $Q_r$ is the
free product of $R_r$ with a tensor algebra on
contractible pairs.

Project to $R_r$ by killing the pairs.
A word in the kernel has a first pair generator,
so it is uniquely $a y_S b$ or $a e_S b$, with
$a$ a word of $R_r$. Set
\[
 h(a y_S b)=(-1)^{|a|}a e_S b,\qquad h(a e_S b)=0.
\]

In $dh+hd$, the terms differentiating $a$ or $b$
cancel, and differentiating this first pair
returns the original word. This contracts the
kernel. Compose the resulting quasi-isomorphism
$Q_r\to R_r$ with $\epsilon$.
\end{proof}

For one-dimensional $V_i$, the squarefree dimensions
make the extra mixed-input data visible:
\[
\begin{array}{c|rrrr}
 r&\dim Q_r^{-3}&\dim Q_r^{-2}&\dim Q_r^{-1}&\dim Q_r^0\\\hline
 2&0&0&2&3\\
 3&0&1&13&13\\
 4&1&32&105&75.
\end{array}
\]
Each row has one degree-zero cohomology class.
The additional generators remember the chosen
paths from mixed inputs to singleton products.

\subsection{Grouping and its higher comparisons}

Let $\mathcal B=(B_1|\cdots|B_m)$ be an ordered
partition of $[r]$, with increasing order inside
each block. Define
\[
 \gamma_{\mathcal B}:R_r\longrightarrow
                           \bigotimes_jR_{B_j}
\]
by sending $q_S$ to the same generator in its
factor if $S\subseteq B_j$, and to zero otherwise.

\begin{proposition}\label{cbc:va:grouping}
The maps $\gamma_{\mathcal B}$ are dg algebra
quasi-isomorphisms and commute with ordered
refinements of the partition.
\end{proposition}
\begin{proof}
For $S$ contained in a block, the chain equation
is unchanged. Otherwise a differential term can
survive only if $S$ meets exactly two blocks and
its two parts are those intersections.
The two orders of that split cancel:
$\operatorname{inv}(L,R)+\operatorname{inv}(R,L)=|L||R|$,
while exchanging their generators contributes
$(-1)^{(1-|L|)(1-|R|)}$.
Their total coefficients are opposite.
This proves the chain equation.

Composing with the tensor of the augmentations
gives $\epsilon$ followed by regrouping the
factors of $P_r$. These augmentations are
quasi-isomorphisms, hence so is $\gamma_{\mathcal B}$.
Under iterated grouping a generator survives
exactly when its colors lie in one final block.
This is also the direct grouping criterion.
\end{proof}

For the relative map, put $T_j=S\cap B_j$.
When $|T_j|\geq2$, let $C_j=c_{T_j}$ and
$E_j=e_{T_j}$. For a singleton use $C_j=x_i$,
$E_j=0$, and for the empty set use $C_j=1$,
$E_j=0$. Let $P_j$ be the increasing singleton
product on $T_j$. Thus $dE_j=C_j-P_j$.

The extension of $\gamma_{\mathcal B}$ to $Q_r$
is
\begin{equation}\label{cbc:va:relative-grouping}
 c_S\longmapsto\prod_jC_j,\qquad
 e_S\longmapsto
       \sum_jC_1\cdots C_{j-1}E_jP_{j+1}\cdots P_m.
\end{equation}
Its differential telescopes to
$\prod_jC_j-\prod_jP_j$.

Ordered refinement expands to the same sum with
one $E$ factor, so these relative maps also
commute with refinement. Their composition with
the augmentations proves they are
quasi-isomorphisms.

Changing the order of two groups changes the
primitive by the boundary
\begin{equation}\label{cbc:va:grouping-square}
 E_1P_2+C_1E_2-P_1E_2-E_1C_2=d(E_1E_2).
\end{equation}
All higher comparisons follow from one closed form.
Let $t_j,dt_j$ have degrees zero and one, and use
the total differential on polynomial forms
tensored with the coefficient algebra. Then
\begin{equation}\label{cbc:va:grouping-cube}
 \Psi=\prod_{j=1}^m
       \bigl(P_j+t_j(C_j-P_j)+dt_jE_j\bigr)
\end{equation}
is closed. Indeed the differential of $t_j(C_j-P_j)$
cancels that of $dt_jE_j$. The factors lie in
separate tensor slots and have total degree zero.

A face sets its corresponding factor to $P_j$
or $C_j$. Integrating along the edges in their
specified order gives
\eqref{cbc:va:relative-grouping}. Square faces
give~\eqref{cbc:va:grouping-square}.

With parameter forms placed first, the top
coefficient is
$(-1)^{m(m-1)/2}E_1\cdots E_m$.
If $I_m$ integrates the top form over the
oriented cube, Stokes' formula gives
\[
 d I_m\Psi=(-1)^{m+1}I_{\partial}\Psi.
\]
All integrals are polynomial over the
characteristic-zero field. The higher comparisons
are therefore finite and retain their specified
faces.

\section{Extension spaces}
\label{cbc:va:extensions}

The resolution turns extension of an action into
equations in its intended target. Their solutions
form a space. Solubility of its first equations,
existence of a full solution, and contractibility
of this space are distinct questions.

Let $E$ be a unital dg algebra and fix
$\phi:F_r\to E$. Write $X_i:V_i\to E^0$ and
$C_S:V_S\to E^0$ for its closed singleton and mixed
images. A map $Q_r\to E$ under $F_r$ consists of
$H_S:V_S\to E^{1-|S|}$ and $T_S:V_S\to E^{-1}$
for $|S|\geq2$, with $H_{\{i\}}=X_i$, satisfying
\begin{align}
 d_EH_S&=
 \sum_{\substack{L\amalg R=S\\L,R\ne\varnothing}}
 (-1)^{|L|+1+\operatorname{inv}(L,R)}H_LH_R,
 \label{cbc:va:extension-h}\\
 d_ET_S&=C_S-X_{i_1}\cdots X_{i_s},
                    \qquad i_1<\cdots<i_s.
 \label{cbc:va:extension-t}
\end{align}
These are exactly the chain equations on the
free generators $q_S,e_S$.

For complexes $U,E$, use the Hom differential
$\partial f=d_Ef-(-1)^{|f|}fd_U$.
The spaces $V_S$ still have degree zero.

\begin{proposition}\label{cbc:va:extension-obstruction}
Fix all $H_T$ for proper nonempty subsets $T\subset S$
with their equations satisfied. The right side
$O_S$ of~\eqref{cbc:va:extension-h} is closed in
$\operatorname{Hom}(V_S,E)^{2-|S|}$.
An extension $H_S$ exists precisely when $[O_S]=0$.
If it exists, its components form a torsor for
$H^{1-|S|}\operatorname{Hom}(V_S,E)$, and at
a chosen solution its higher homotopy groups are
\[
 \pi_j=H^{1-|S|-j}\operatorname{Hom}(V_S,E),
                         \qquad j\geq1.
\]
All lower-stage maps are fixed in this assertion.
For a mixed cell, the obstruction lies in
$H^0\operatorname{Hom}(V_S,E)$ and its choice
groups are $H^{-1-j}\operatorname{Hom}(V_S,E)$
for $j\geq0$.
\end{proposition}
\begin{proof}
Differentiating $O_S$ produces the two opposite
terms for each ordered three-part decomposition
of $S$, as in $d^2q_S=0$. Thus it is closed.
The next equation is $\partial H_S=O_S$, which
has a solution exactly when this class vanishes.

Subtract a chosen solution. The fiber becomes the
simplicial abelian group of closed degree
$1-|S|$ maps from $V_S$ into
$E\otimes\Omega^\bullet(\Delta^\bullet)$.
Integration of polynomial forms identifies its
normalized homotopy groups with the Hom
cohomology in degrees $1-|S|-j$.
In particular, an interval changes two solutions
by a boundary, giving the stated torsor.
The same argument applies to a generator of
degree minus one.
\end{proof}

A contraction $h$ of the relevant Hom complex
gives $H_S=hO_S$ when its cohomology projection
vanishes. If units or collision values are fixed,
this must be a contraction of the corresponding
relative Hom complex. A contraction that moves
those values does not solve the extension problem.

There is also a model of the complete solution
space:
\begin{equation}\label{cbc:va:full-extension-fibre}
 \operatorname{Map}_{F_r/}(Q_r,E)
 \simeq
 \operatorname{fib}_{\phi}
 \left(
 \operatorname{Map}_{\mathrm{dgAlg}}(P_r,E)
 \longrightarrow
 \operatorname{Map}_{\mathrm{dgAlg}}(F_r,E)
 \right).
\end{equation}
Indeed $F_r$ is free, $F_r\to Q_r$ is a
cofibration, and $Q_r\to P_r$ is a
quasi-isomorphism. Polynomial forms compute
mapping spaces from these cofibrant sources.
Restriction is a Kan fibration, so its strict
fiber computes its homotopy fiber. These are
the unital associative instances of
\cite{Hinich1997}*{Theorem~4.1.1 and \S\S4.8.1--4.8.8}.

If $E$ is zero below degree $-N$, with $N\geq0$,
only $|S|\leq N+2$ can give nonzero higher-cell
obstructions. At the last of these sizes the
obstruction can survive although its primitive
would have to lie in a zero degree. Larger
right sides vanish as elements of $E$.
A cohomological bound alone does not give this
strict bound on the equations.

\subsection{Choices that survive and choices that disappear}

Take $r=2$, $V_1=V_2=k$, and
$E=k\oplus k\eta$ with $|\eta|=-1$, $d=0$,
$\eta^2=0$. Prescribe $X_1=X_2=C_{12}=0$.
Then
\[
 H_{12}=\lambda\eta,\qquad T_{12}=\mu\eta
                         \quad(\lambda,\mu\in k)
\]
are all the solutions.

A degree-minus-one
element of $E\otimes\Omega^\bullet(\Delta^n)$
is $\eta f$, with $f$ a polynomial function.
Closedness forces $df=0$, hence $f$ is constant.
Thus~\eqref{cbc:va:full-extension-fibre} is the
discrete space $k^2$. Existence does not imply
contractibility.

An intermediate choice need not survive a later
cell. With the same $E$, maps from
$k\langle x\rangle$, $|x|=-1$, $dx=0$, have
components $x\mapsto\lambda\eta$.
Adjoin $y$ in degree minus two with $dy=x$.
In every polynomial-form simplex, the image of
$y$ is zero for degree reasons. Its differential
then forces the image of $x$ to be zero.

The complete extension space is a point.
Consequently a stage choice group alone cannot
disprove contractibility of the full space.

\subsection{A triple obstruction after all pair choices}

Start with the free graded algebra on $x_1,x_2,x_3$
in degree zero and $h_{12},h_{13},h_{23}$ in
degree minus one, with
$dh_{ij}=x_ix_j-x_jx_i$.
Give $x_i$ color weight $e_i$ and $h_{ij}$ color
weight $e_i+e_j$. Quotient by every word in
which a color occurs twice, and call the result
$E$. The differential preserves color counts,
so the quotient is a dg algebra.

Its degree-zero basis consists of the unit,
three singletons, six ordered pairs and six
ordered triples. Its degree-minus-one basis
consists of the three $h_{ij}$ and the six
products of an $h_{ij}$ with its complementary
singleton, on either side. Thus $\dim E=25$.
Define
\begin{equation}\label{cbc:va:omega}
 \omega=x_1h_{23}-x_2h_{13}+x_3h_{12}
             -h_{12}x_3+h_{13}x_2-h_{23}x_1.
\end{equation}

\begin{theorem}\label{cbc:va:triple-obstruction}
Every pair of the marked maps $k[x_i]\to E$
extends to the corresponding resolution $R_2$.
No choice of pair primitives extends them to
$R_3$. More precisely,
\[
 H^0(E)=k[x_1,x_2,x_3]/(x_1^2,x_2^2,x_3^2),
 \qquad H^{-1}(E)=k\omega,
\]
and every choice gives the triple obstruction
$\omega$.
\end{theorem}
\begin{proof}
The displayed $h_{ij}$ solve the pair equations.
In each pair weight, their differentials are
injective. In weight $(1,1,1)$, the six
degree-minus-one basis vectors map to the six
ordered triple words by the incidence matrix
of a hexagon. Its rank is five and its kernel
is its oriented cycle~\eqref{cbc:va:omega}.

The total differential therefore has rank eight.
Its degree-zero quotient identifies the orderings
of each set of distinct colors and kills every
repeated-color product. This gives the two
cohomology formulas.

Every other pair primitive is
$h_{ij}+\lambda_{ij}\omega$.
The augmentation ideal annihilates $\omega$
on both sides, since its terms already use all
three colors. Substituting any such primitives
in~\eqref{cbc:va:triple-differential} therefore
leaves $\omega$ unchanged.
Its six summands are distinct basis vectors,
so it is nonzero. Since $E^{-2}=0$, it cannot
be a boundary.
\end{proof}

Prescribe $c_S$ to be the increasing singleton
product. Then $T_S=0$ solves all mixed equations,
and the same obstruction excludes a map
$Q_3\to E$ under these data.
The regular action of $E$ detects its class:
multiplication by $\omega$ sends $[1]$ to
$[\omega]\ne0$. This detects the specified lifts
in $\operatorname{End}_k(E)$. It does not
classify the additional pair primitives in that
larger algebra.

The obstruction also gives an integral family.
Over $k[s]$, adjoin a vector $K$ of degree
minus two with $dK=s\omega$.
Set $K^2=0$ and let the augmentation ideal
annihilate $K$ on both sides.
These relations are preserved by the differential
because the same ideal annihilates $\omega$.
The resulting $E_s$ is graded-free of rank $26$.

\begin{proposition}\label{cbc:va:parameter-obstruction}
One has
\[
 H^0(E_s)=k[s]\otimes H^0(E),\qquad
 H^{-1}(E_s)=k[s]/(s)\cdot\omega,\qquad H^{-2}(E_s)=0.
\]
Let $R$ be a commutative $k[s]$-algebra.
The prescribed singleton maps extend to $R_3$ after
base change to $R$ exactly when $s$ is a unit there. This criterion allows every choice
of pair primitives. It also applies to $Q_3$
when the mixed inputs are their prescribed
singleton products.
\end{proposition}
\begin{proof}
The differential in degrees minus one and zero
is split over $k$, with kernel $k\omega$.
After base change its kernel is $R\omega$,
so changes of pair primitives still leave the
triple expression unchanged.
An image of $q_{123}$ must be $bK$ and solve
$bs\omega=\omega$. As $\omega$ is part of a
free $R$-basis, this is $bs=1$.

Over $k[s]$, the remaining negative complex is
$k[s]K\xrightarrow{s}k[s]\omega$.
Its kernel is zero and its cokernel is
$k[s]/(s)\cdot\omega$. The degree-zero
differential is unchanged.
\end{proof}

The primitive over $k[s,s^{-1}]$ is $K/s$.
There is none over $k[s]$ or $k[[s]]$.
At $s=0$, both $K$ and $\omega$ survive.
At a nonzero field value, rescaling $K$
gives the quotient of $R_3$ by repeated-color
words. Graded flatness of the family therefore
does not imply flatness of its cohomology.

More generally, over a commutative $k$-algebra
$R$, adjoin vectors $K_j$ with
$dK_j=f_j\omega$, zero mutual products and
zero products with the augmentation ideal.
The same calculation gives
\[
 H^{-1}=R/(f_1,\ldots,f_m)\cdot\omega,\qquad
 H^{-2}=\ker\bigl(R^m\xrightarrow{(f_1,\ldots,f_m)}R\bigr).
\]
A triple extension exists exactly when
$1\in(f_1,\ldots,f_m)$. The degree-minus-two
classes are the syzygies among these coefficients.

\section{The bar of the full action diagram}
\label{cbc:va:action-bar}

An extension can become possible after mapping
to another acting object. To decide a universal
property, one must therefore retain the arrows
between all permitted actions and their comparison
paths. The bar of that diagram gives an explicit
chain model for the resulting coefficient map.

First let $\mathcal A$ be a small ordinary
category and $X:\mathcal A\to\operatorname{Ch}(k)$
a covariant diagram. Its simplicial replacement
has the length-$n$ term
\[
 \bigoplus_{a_0\xrightarrow{f_1}\cdots
                    \xrightarrow{f_n}a_n}X(a_0)[n].
\]

Use the direct-sum total complex $\mathcal B(X)$.
An element $[f_1|\cdots|f_n]x$, with $|x|=q$
internally, has total degree $q-n$.
The first face applies $X(f_1)$, the interior
faces compose adjacent arrows, and the last
face deletes $f_n$. Degeneracies insert identities.
If $\partial$ is the alternating face differential,
put $d=\partial+(-1)^nd_X$ on length $n$.

For a second diagram $Y$, define the
Alexander--Whitney map by
\begin{equation}\label{cbc:va:aw}
\begin{aligned}
 \operatorname{AW}([f_1|\cdots|f_n](x\otimes y))
  =\sum_{p=0}^n(-1)^{(n-p)|x|}
    [f_1|\cdots|f_p]x\otimes
    [f_{p+1}|\cdots|f_n]Y(f_p\cdots f_1)y.
\end{aligned}
\end{equation}
The empty composite for $p=0$ is the identity.
The second string begins at $a_p$. Its coefficient
must be transported to $Y(a_p)$.
The sign moves $x$ past its suffix, of degree
$-(n-p)$.

\begin{proposition}\label{cbc:va:aw-chain}
Formula~\eqref{cbc:va:aw} defines a natural
coassociative chain map
\[
 \mathcal B(X\otimes Y)\longrightarrow
                         \mathcal B(X)\otimes\mathcal B(Y).
\]
It descends to normalization and represents the
map of homotopy colimits induced by
$\mathcal A\to\mathcal A^2$, $a\mapsto(a,a)$.
\end{proposition}
\begin{proof}
Before moving coefficients through strings,
the formula splits a simplex into a prefix
and a suffix. Deleting the prefix vertices
from the second factor applies its first
face repeatedly, giving $Y(f_p\cdots f_1)y$.

A face strictly inside either interval agrees
with that face taken before the split.
At the splitting vertex, the last face of
one prefix cancels the first face of the
next suffix. At the initial vertex,
functoriality identifies the composite
coefficient transport with its first face.
The last vertex requires no transport.
These account for the face differential.

Restoring the displayed coefficient order gives
$(-1)^{(n-p)|x|}$. The same permutation carries
the internal differential: its source
simplex sign is $(-1)^n$, and the first
output has total degree $|x|-p$.
Thus it also satisfies the internal chain
equation. An identity arrow remains in a
prefix or a suffix in every summand, which
proves descent to normalization.

For three diagrams, either iterated diagonal
sums over $0\leq p\leq q\leq n$.
Both transport the second and third coefficients
to $a_p,a_q$ and have sign
$(-1)^{(n-p)|x|+(n-q)|y|}$.
This proves coassociativity.

The diagonal category functor first sends an
arrow string to two identical strings.
The prefix--suffix map compares the diagonal
of their bisimplicial replacement with its
total complex. The shuffle map gives the
reverse comparison, with the degeneracy
homotopies. The latter total complex is
$\mathcal B(X)\otimes\mathcal B(Y)$.
Hence the map has the stated derived meaning.
\end{proof}

For one arrow $f:a\to b$, the formula reads
\[
 \operatorname{AW}([f](x\otimes y))
 =(-1)^{|x|}[\,]_ax\otimes[f]y
                   +[f]x\otimes[\,]_bY(f)y.
\]
Both the transport and the sign are necessary.

Put $Q(a,b)=X(a)\otimes Y(b)$. Over the field $k$,
tensor product preserves colimits and computes
derived tensor product. The fiber of the
diagonal-induced comparison therefore has model
\begin{align}
 \mathcal K^j&=\mathcal B(X\otimes Y)^j
     \oplus(\mathcal B(X)\otimes\mathcal B(Y))^{j-1},
 \notag\\
 d(u,v)&=(du,\operatorname{AW}(u)-dv).
 \label{cbc:va:action-bar-fibre}
\end{align}
This describes its differential, without
asserting that the fiber is acyclic.
Over a general coefficient ring one first
takes derived-flat replacements.

For an action category with simplicial mapping
spaces, use its coherent simplicial replacement
and normalize every simplicial direction.
The diagonal and shuffle maps in those
directions give the same derived comparison.
Passing to the homotopy category would omit
the action paths that enter this replacement.

On a binary block-disjoint open, suppose each
factorizing actor's coefficient is the external product
of its two block coefficients.
Require coherence over the full action index. Their diagram comparison then has the
form~\eqref{cbc:va:action-bar-fibre}, with external
products. Universal factorization requires its
acyclicity on that full index. A class on a
restricted index may disappear after adding
another recipient and its comparison simplices.

Even bar length matters. For the one-object
identity category with constant coefficient
$k$, the unnormalized bar has one copy of
$k$ in every length and boundary
$\partial_n=\sum_{i=0}^n(-1)^i$.
Truncating after length one creates
$H^{-1}=k$ because $\partial_1=0$.
In the full bar, $\partial_2=1$ kills this
class. Normalization removes all positive
lengths. A finite unnormalized truncation
can therefore create an obstruction absent
from the full diagram.

The remaining geometric comparison must join
these separated relations to the relations
on singleton collisions. It must also retain
the positive local vacuum and the specified
action paths on the full ordered atlas.
For the dual-number pointed actor on $C_{V_\kappa}$,
the contractibility question requires these cells
in its actual relative endomorphism complex. A nonzero obstruction
must survive every permitted lower-cell
choice there, and a nontrivial choice must
survive the rest of the extension tower.
The ordinary coefficient algebras above
specify these tests. They do not themselves
construct the required supported comparison.
